
\documentclass[11pt,oneside]{book}

\usepackage{iftex}
\ifPDFTeX
  \usepackage[T1]{fontenc}
  \usepackage[utf8]{inputenc}
  \usepackage{textcomp}
  \usepackage{lmodern}
  \usepackage{amsmath,amssymb,amsthm,mathtools}
\else
  \usepackage{fontspec}
  \usepackage{unicode-math}
  \defaultfontfeatures{Scale=MatchLowercase}
  \defaultfontfeatures[\rmfamily]{Ligatures=TeX,Scale=1}
  \setmainfont{DejaVu Serif}
  \setsansfont{DejaVu Sans}
  \setmonofont{DejaVu Sans Mono}
  \usepackage{amsmath,amsthm,mathtools}
\fi
\usepackage{microtype}
\usepackage[margin=1in]{geometry}
\usepackage{xcolor}
\usepackage{fancyvrb}
\usepackage{longtable,booktabs,array,calc}
\usepackage{etoolbox}
\usepackage{enumitem}
\usepackage{bookmark}
\usepackage{hyperref}
\IfFileExists{xurl.sty}{\usepackage{xurl}}{}
\urlstyle{same}
\hypersetup{hidelinks,pdfcreator={ChatGPT}}

\newcommand{\VerbBar}{|}
\newcommand{\VERB}{\Verb[commandchars=\\\{\}]}
\DefineVerbatimEnvironment{Highlighting}{Verbatim}{commandchars=\\\{\}}
\newenvironment{Shaded}{}{}
\newcommand{\AlertTok}[1]{\textcolor[rgb]{1.00,0.00,0.00}{\textbf{#1}}}
\newcommand{\AnnotationTok}[1]{\textcolor[rgb]{0.38,0.63,0.69}{\textbf{\textit{#1}}}}
\newcommand{\AttributeTok}[1]{\textcolor[rgb]{0.49,0.56,0.16}{#1}}
\newcommand{\BaseNTok}[1]{\textcolor[rgb]{0.25,0.63,0.44}{#1}}
\newcommand{\BuiltInTok}[1]{\textcolor[rgb]{0.00,0.50,0.00}{#1}}
\newcommand{\CharTok}[1]{\textcolor[rgb]{0.25,0.44,0.63}{#1}}
\newcommand{\CommentTok}[1]{\textcolor[rgb]{0.38,0.63,0.69}{\textit{#1}}}
\newcommand{\CommentVarTok}[1]{\textcolor[rgb]{0.38,0.63,0.69}{\textbf{\textit{#1}}}}
\newcommand{\ConstantTok}[1]{\textcolor[rgb]{0.53,0.00,0.00}{#1}}
\newcommand{\ControlFlowTok}[1]{\textcolor[rgb]{0.00,0.44,0.13}{\textbf{#1}}}
\newcommand{\DataTypeTok}[1]{\textcolor[rgb]{0.56,0.13,0.00}{#1}}
\newcommand{\DecValTok}[1]{\textcolor[rgb]{0.25,0.63,0.44}{#1}}
\newcommand{\DocumentationTok}[1]{\textcolor[rgb]{0.73,0.13,0.13}{\textit{#1}}}
\newcommand{\ErrorTok}[1]{\textcolor[rgb]{1.00,0.00,0.00}{\textbf{#1}}}
\newcommand{\ExtensionTok}[1]{#1}
\newcommand{\FloatTok}[1]{\textcolor[rgb]{0.25,0.63,0.44}{#1}}
\newcommand{\FunctionTok}[1]{\textcolor[rgb]{0.02,0.16,0.49}{#1}}
\newcommand{\ImportTok}[1]{\textcolor[rgb]{0.00,0.50,0.00}{\textbf{#1}}}
\newcommand{\InformationTok}[1]{\textcolor[rgb]{0.38,0.63,0.69}{\textbf{\textit{#1}}}}
\newcommand{\KeywordTok}[1]{\textcolor[rgb]{0.00,0.44,0.13}{\textbf{#1}}}
\newcommand{\NormalTok}[1]{#1}
\newcommand{\OperatorTok}[1]{\textcolor[rgb]{0.40,0.40,0.40}{#1}}
\newcommand{\OtherTok}[1]{\textcolor[rgb]{0.00,0.44,0.13}{#1}}
\newcommand{\PreprocessorTok}[1]{\textcolor[rgb]{0.74,0.48,0.00}{#1}}
\newcommand{\RegionMarkerTok}[1]{#1}
\newcommand{\SpecialCharTok}[1]{\textcolor[rgb]{0.25,0.44,0.63}{#1}}
\newcommand{\SpecialStringTok}[1]{\textcolor[rgb]{0.73,0.40,0.53}{#1}}
\newcommand{\StringTok}[1]{\textcolor[rgb]{0.25,0.44,0.63}{#1}}
\newcommand{\VariableTok}[1]{\textcolor[rgb]{0.10,0.09,0.49}{#1}}
\newcommand{\VerbatimStringTok}[1]{\textcolor[rgb]{0.25,0.44,0.63}{#1}}
\newcommand{\WarningTok}[1]{\textcolor[rgb]{0.38,0.63,0.69}{\textbf{\textit{#1}}}}
\providecommand{\tightlist}{\setlength{\itemsep}{0pt}\setlength{\parskip}{0pt}}
\setlength{\emergencystretch}{3em}
\makeatletter
\patchcmd\longtable{\par}{\if@noskipsec\mbox{}\fi\par}{}{}
\makeatother
\IfFileExists{footnotehyper.sty}{\usepackage{footnotehyper}}{\usepackage{footnote}}
\makesavenoteenv{longtable}

\theoremstyle{plain}
\providecommand{\blacksquare}{\rule{1.1ex}{1.1ex}}
\newtheorem{theorem}{Theorem}[section]
\newtheorem{proposition}[theorem]{Proposition}
\newtheorem{lemma}[theorem]{Lemma}
\newtheorem{corollary}[theorem]{Corollary}
\theoremstyle{definition}
\newtheorem{definition}[theorem]{Definition}
\newtheorem{example}[theorem]{Example}
\newtheorem{exercise}[theorem]{Exercise}
\theoremstyle{remark}
\newtheorem{remark}[theorem]{Remark}

\title{Cyclic Codes}
\author{}
\date{}

\begin{document}
\frontmatter
\maketitle
\tableofcontents
\mainmatter


\chapter{Cyclic Codes}\label{chap:cyclic-codes}

This chapter presents a systematic introduction to cyclic codes. The exposition emphasizes the algebraic structure of cyclic codes, their concrete realization by generator and parity-check polynomials, their implementation by shift-register circuits, and their role in both classical and modern applications.



\section{Historical Background}\label{sec:historical-background}

We briefly record several principal milestones in the development of cyclic codes and their surrounding theory.

\renewcommand{\arraystretch}{1.1}
\begin{longtable}{@{}p{1.8cm}p{4.2cm}>{\raggedright\arraybackslash}p{8.0cm}@{}}
\toprule
\textbf{Year} & \textbf{Milestone} & \textbf{Description} \\
\midrule
\endfirsthead

\toprule
\textbf{Year} & \textbf{Milestone} & \textbf{Description} \\
\midrule
\endhead

\bottomrule
\endfoot

1948
& Foundation of coding theory
& Shannon established information theory and proved the existence of reliable codes for noisy channels {[}1{]}. \\

1950
& Hamming codes
& Hamming introduced the first systematic linear codes capable of correcting a single error {[}2{]}. \\

1954
& Reed--Muller codes
& Reed and Muller introduced a family of multiple-error-correcting codes arising from Boolean functions {[}3{]}{[}4{]}. \\

1957
& Algebraic formulation of cyclic codes
& Prange showed that cyclic codes may be identified with ideals in the quotient ring \(\mathbb{F}_q[X]/\langle X^n-1\rangle\) {[}5{]}. \\

1959--1960
& BCH codes
& Hocquenghem, Bose, and Ray-Chaudhuri introduced BCH codes, thereby providing a systematic algebraic construction with prescribed designed distance {[}6{]}{[}7{]}. \\

1960
& Reed--Solomon codes
& Reed and Solomon introduced an important class of nonbinary cyclic codes with strong symbol-level correction capability {[}8{]}. \\

1961
& Algebraic decoding
& Peterson developed an algebraic decoding method for BCH codes and published a foundational text on error-correcting codes {[}9{]}. \\

1968
& Berlekamp--Massey algorithm
& Berlekamp and Massey introduced an efficient iterative procedure for determining the error-locator polynomial {[}10{]}{[}11{]}. \\

1970s
& Standardization of CRC
& Cyclic redundancy check became a standard method of error detection because cyclic codes admit efficient LFSR-based implementation. \\

1977
& Goppa codes and code-based cryptography
& Goppa introduced algebraic geometry codes, and McEliece employed Goppa codes in a foundational public-key cryptosystem {[}12{]}{[}13{]}. \\

Since the 1990s
& Modern developments
& Despite the emergence of turbo codes and LDPC codes, cyclic codes have remained important because of their algebraic transparency, implementation efficiency, and cryptographic relevance {[}14{]}{[}15{]}. \\

\end{longtable}

\subsection{Change of perspective}\label{sec:change-of-perspective}

The historical development of cyclic codes also reflects a change in mathematical perspective. In the early stage, practical engineering demands led to the discovery of algebraic structure. In the middle stage, that algebraic structure in turn guided the systematic construction of codes and decoding algorithms. In the modern stage, cyclic and quasi-cyclic structures continue to interact with hardware design, complexity theory, and post-quantum cryptography. 

\section{Definition of Cyclic Codes}\label{sec:definition-of-cyclic-codes}

A cyclic code is a linear block code closed under cyclic shift. More precisely, a linear code \(\mathcal{C}\) of length \(n\) is called \emph{cyclic} if it satisfies the following two conditions.

\begin{enumerate}
\def\labelenumi{\arabic{enumi}.}
\tightlist
\item
  \textbf{Linearity}: for any two codewords
  \(\mathbf{c}_1, \mathbf{c}_2 \in \mathcal{C}\), we have
  \(\mathbf{c}_1 + \mathbf{c}_2 \in \mathcal{C}\).
\item
  \textbf{Cyclic shift property}: if \(\mathbf{c}=(c_0,c_1,\dots,c_{n-1})\in\mathcal{C}\), then its one-step cyclic shift
  \[\rho(\mathbf{c})=(c_{n-1},c_0,c_1,\dots,c_{n-2})\]
  also belongs to \(\mathcal{C}\).
\end{enumerate}

Therefore any cyclic shift by two steps, three steps, or in fact any
number of steps also yields a valid codeword.

\subsubsection{Example 1.1: Checking a Cyclic
Shift}\label{example-1.1-checking-a-cyclic-shift}

In a binary \([7,4]_2\) cyclic code, suppose
\[\mathbf{c} = (1,0,0,1,0,1,1)\] is a codeword. Then:

\begin{itemize}
\tightlist
\item
  1-step shift: \(\rho(\mathbf{c})=(1,1,0,0,1,0,1)\)
  $\to$ codeword
\item
  2-step shift: \(\rho^2(\mathbf{c})=(1,1,1,0,0,1,0)\)
  $\to$ codeword
\item
  3-step shift: \(\rho^3(\mathbf{c})=(0,1,1,1,0,0,1)\)
  $\to$ codeword
\end{itemize}

All of these vectors must belong to the code space \(\mathcal{C}\).

\subsubsection{Example 1.2: Deciding Whether a Code Is
Cyclic}\label{example-1.2-deciding-whether-a-code-is-cyclic}

Let \(n=4\) and consider the set
\[\mathcal{C} = \{(0,0,0,0),\ (1,0,1,0),\ (0,1,0,1),\ (1,1,1,1)\}.\]

\begin{itemize}
\tightlist
\item
  Linearity check:
  \((1,0,1,0) \oplus (0,1,0,1) = (1,1,1,1) \in \mathcal{C}\).
\item
  Cyclic-shift check:

  \begin{itemize}
  \tightlist
  \item
    \((1,0,1,0)\) shifted once gives \((0,1,0,1) \in \mathcal{C}\).
  \item
    \((0,1,0,1)\) shifted once gives \((1,0,1,0) \in \mathcal{C}\).
  \item
    \((1,1,1,1)\) shifted once gives \((1,1,1,1) \in \mathcal{C}\).
  \end{itemize}
\end{itemize}

Hence this is a cyclic code.

\subsubsection{Example 1.3: A Linear Code That Is Not
Cyclic}\label{example-1.3-a-linear-code-that-is-not-cyclic}

Let
\[\mathcal{C}' = \{(0,0,0,0),\ (1,1,0,0),\ (0,0,1,1),\ (1,1,1,1)\}.\]

Shifting \((1,1,0,0)\) one step cyclically gives \((0,1,1,0)\), which is
not in \(\mathcal{C}'\). Therefore \(\mathcal{C}'\) is linear, but
\textbf{not} cyclic.


\section{Implementation Advantage of Cyclic-Shift Invariance}\label{sec:implementation-advantage}

Cyclic codes were originally studied not because of abstract algebra,
but because of \textbf{hardware efficiency}. Prange's key observation in
the 1950s was that if a code is closed under cyclic shift, then encoding
and error detection can be implemented using nothing more than shift
registers.

\subsubsection{The Drawback of General Linear
Codes}\label{the-drawback-of-general-linear-codes}

For a general linear \([n,k]_q\) code, encoding is done by matrix
multiplication: \[\mathbf{c} = \mathbf{m} \cdot G.\] This means:

\begin{itemize}
\tightlist
\item
  the \(k \times n\) generator matrix \(G\) must be stored in memory,
\item
  \(k \cdot n\) additions and multiplications are needed,
\item
  computation cannot begin until the whole message block has arrived.
\end{itemize}

For example, for a \([31,21]_2\) code, one must store
\(21 \times 31 = 651\) matrix entries and perform 651 binary operations
for every block.

\subsubsection{The Cyclic-Code Alternative: Shift
Registers}\label{the-cyclic-code-alternative-shift-registers}

For cyclic codes, every codeword is a multiple of the generator
polynomial \(g(X)\), and encoding reduces to polynomial multiplication
or polynomial division with remainder. These operations can be
implemented using an LFSR (linear feedback shift register):

\begin{itemize}
\tightlist
\item
  only \(r=n-k\) one-bit registers are required,
\item
  the input stream is processed one bit per clock,
\item
  after the last bit passes through the circuit, the register contents
  are exactly the parity bits or syndrome.
\end{itemize}

For a \([31,21]_2\) code, 10 registers plus a few XOR gates are enough.

\subsubsection{\texorpdfstring{Example: Matrix Multiplication vs.~LFSR
for a \([7,4]_2\)
Code}{Example: Matrix Multiplication vs.~LFSR for a {[}7,4{]}\_2 Code}}\label{example-matrix-multiplication-vs.-lfsr-for-a-74_2-code}

\begin{longtable}[]{@{}
  >{\raggedright\arraybackslash}p{(\columnwidth - 4\tabcolsep) * \real{0.3333}}
  >{\raggedright\arraybackslash}p{(\columnwidth - 4\tabcolsep) * \real{0.3333}}
  >{\raggedright\arraybackslash}p{(\columnwidth - 4\tabcolsep) * \real{0.3333}}@{}}
\toprule\noalign{}
\begin{minipage}[b]{\linewidth}\raggedright
Item
\end{minipage} & \begin{minipage}[b]{\linewidth}\raggedright
General linear code (matrix multiplication)
\end{minipage} & \begin{minipage}[b]{\linewidth}\raggedright
Cyclic code (LFSR)
\end{minipage} \\
\midrule\noalign{}
\endhead
\bottomrule\noalign{}
\endlastfoot
Storage & \(G \in GF(2)^{4 \times 7}\) $\to$ 28 bits &
\(g(X)=1+X+X^3\) $\to$ 4 bits (coefficients \([1,1,0,1]\)) \\
Operations per block & \(4 \times 7 = 28\) AND/XOR operations & 7 clocks
(roughly one XOR per bit) \\
Processing mode & Batch processing after full block arrival & Real-time,
bit-by-bit \\
Latency & Proportional to block length & Almost zero \\
Circuit complexity & \(n\) XOR trees driven by \(k\) inputs & \(r\)
flip-flops plus XOR gates \\
\end{longtable}

\subsubsection{Detailed Analysis: Why Does It Take 7 Clocks?
(Non-Systematic
Encoding)}\label{detailed-analysis-why-does-it-take-7-clocks-non-systematic-encoding}

Non-systematic encoding \(c(X)=m(X)g(X)\) is, mathematically, a
convolution of two polynomials. For a \([7,4]_2\) code the message has 4
bits, the generator polynomial has degree 3, and the codeword has length
7.

\paragraph{1. Coefficient Expansion of the Polynomial
Product}\label{coefficient-expansion-of-the-polynomial-product}

Let \[m(X)=m_0+m_1X+m_2X^2+m_3X^3,\] \[g(X)=g_0+g_1X+g_2X^2+g_3X^3.\]
Then in \[c(X)=\sum_{k=0}^{6} c_k X^k,\] we have:

\begin{itemize}
\tightlist
\item
  \(c_0 = m_0 g_0\)
\item
  \(c_1 = m_0 g_1 + m_1 g_0\)
\item
  \(c_2 = m_0 g_2 + m_1 g_1 + m_2 g_0\)
\item
  \(c_3 = m_0 g_3 + m_1 g_2 + m_2 g_1 + m_3 g_0\)
\item
  \(c_4 = m_1 g_3 + m_2 g_2 + m_3 g_1\)
\item
  \(c_5 = m_2 g_3 + m_3 g_2\)
\item
  \(c_6 = m_3 g_3\)
\end{itemize}

\paragraph{2. Clock-by-Clock Operation (Serial
Processing)}\label{clock-by-clock-operation-serial-processing}

In a serial implementation, one message bit enters per clock, and one
codeword bit is emitted per clock.

\begin{longtable}[]{@{}
  >{\centering\arraybackslash}p{(\columnwidth - 6\tabcolsep) * \real{0.2778}}
  >{\centering\arraybackslash}p{(\columnwidth - 6\tabcolsep) * \real{0.2778}}
  >{\centering\arraybackslash}p{(\columnwidth - 6\tabcolsep) * \real{0.2778}}
  >{\raggedright\arraybackslash}p{(\columnwidth - 6\tabcolsep) * \real{0.1667}}@{}}
\toprule\noalign{}
\begin{minipage}[b]{\linewidth}\centering
Clock
\end{minipage} & \begin{minipage}[b]{\linewidth}\centering
Input
\end{minipage} & \begin{minipage}[b]{\linewidth}\centering
Output
\end{minipage} & \begin{minipage}[b]{\linewidth}\raggedright
Explanation
\end{minipage} \\
\midrule\noalign{}
\endhead
\bottomrule\noalign{}
\endlastfoot
1 & \(m_0\) & \(c_0\) & The first codeword bit appears immediately \\
2 & \(m_1\) & \(c_1\) & Contributions of \(m_0\) and \(m_1\) combine \\
3 & \(m_2\) & \(c_2\) & \\
4 & \(m_3\) & \(c_3\) & Message input finishes, but the tail has not
flushed out yet \\
5 & \(0\) & \(c_4\) & Feed zeros to flush the remaining state \\
6 & \(0\) & \(c_5\) & \\
7 & \(0\) & \(c_6\) & Final output bit \\
\end{longtable}

\paragraph{3. Hardware Structure (ASCII
Art)}\label{hardware-structure-ascii-art}

For \(g(X)=1+X+X^3\):

\begin{Shaded}
\begin{Highlighting}[]
\NormalTok{       Message input m\_k}
\NormalTok{        |}
\NormalTok{        +{-}{-}{-}{-}{-}{-}{-}{-}{-}{-}{-}{-}{-}{-}{-}+{-}{-}{-}{-}{-}{-}{-}{-}{-}{-}{-}{-}{-}{-}{-}+{-}{-}{-}{-}{-}{-}{-}{-}{-}{-}{-}{-}{-}{-}{-}+}
\NormalTok{        |               |               |               |}
\NormalTok{       (x) g0=1        (x) g1=1        (x) g2=0        (x) g3=1}
\NormalTok{        |               |               |               |}
\NormalTok{        v               v               v               v}
\NormalTok{ 0 {-}{-}\textgreater{}(+){-}{-}{-}{-}[ D0 ]{-}{-}{-}{-}(+){-}{-}{-}{-}[ D1 ]{-}{-}{-}{-}(+){-}{-}{-}{-}[ D2 ]{-}{-}{-}{-}(+){-}{-}{-}{-}\textgreater{} Codeword output c\_k}
\end{Highlighting}
\end{Shaded}

\textbf{How it works:} 1. \textbf{Input phase (clocks 1-4):} message
bits \(m_0,m_1,m_2,m_3\) are fed in sequence. At each clock the current
input bit is multiplied by each coefficient of \(g(X)\) and combined
with delayed values stored in the registers. 2. \textbf{Flush phase
(clocks 5-7):} after the message bits are exhausted, zeros are shifted
in. The intermediate values stored in \(D_0,D_1,D_2\) are flushed out,
yielding \(c_4,c_5,c_6\). 3. \textbf{Result:} after 7 clocks all
registers are cleared and the 7-bit codeword has been fully emitted.

Hence the total number of clocks is \[4 + 3 = 7,\] which equals the
codeword length \(n\).

\paragraph{4. Step-by-Step
Visualization}\label{step-by-step-visualization}

Let the message be \(\mathbf{m}=(1,1,0,1)\) and let \(g(X)=1+X+X^3\).

\textbf{Initial state:} all registers are zero.

\begin{Shaded}
\begin{Highlighting}[]
\NormalTok{Input m\_k: {-}}
\NormalTok{Registers: [ D0:0 ] [ D1:0 ] [ D2:0 ]}
\NormalTok{Output c\_k: {-}}
\end{Highlighting}
\end{Shaded}

\textbf{Clock 1:} input \(m_0=1\)

\begin{itemize}
\tightlist
\item
  \(c_0 = 1 \cdot 1 \oplus 0 = 1\)
\item
  New state: \(D_0 \leftarrow 1\), \(D_1 \leftarrow 0\),
  \(D_2 \leftarrow 1\)
\end{itemize}

\begin{Shaded}
\begin{Highlighting}[]
\NormalTok{Input m\_k: 1}
\NormalTok{     |}
\NormalTok{     v (x1)   (x1)   (x0)   (x1)}
\NormalTok{     |   |      |      |      |}
\NormalTok{0 {-}{-}(+){-}{-}[ D0 ]{-}{-}(+){-}{-}[ D1 ]{-}{-}(+){-}{-}[ D2 ]{-}{-}(+){-}{-}\textgreater{} Output c\_0: 1}
\NormalTok{         (New:1)      (New:0)      (New:1)}
\end{Highlighting}
\end{Shaded}

\textbf{Clock 2:} input \(m_1=1\)

\begin{itemize}
\tightlist
\item
  \(c_1 = 1 \cdot 1 \oplus D_0 = 1 \oplus 1 = 0\)
\item
  New state: \(D_0 \leftarrow 1\), \(D_1 \leftarrow 1\),
  \(D_2 \leftarrow 1\)
\end{itemize}

\begin{Shaded}
\begin{Highlighting}[]
\NormalTok{Input m\_k: 1}
\NormalTok{     |}
\NormalTok{     v}
\NormalTok{0 {-}{-}(+){-}{-}[ D0:1 ]{-}{-}(+){-}{-}[ D1:0 ]{-}{-}(+){-}{-}[ D2:1 ]{-}{-}(+){-}{-}\textgreater{} Output c\_1: 0}
\NormalTok{         (New:1)        (New:1)        (New:1)}
\end{Highlighting}
\end{Shaded}

\textbf{Clock 3:} input \(m_2=0\)

\begin{itemize}
\tightlist
\item
  \(c_2 = 0 \cdot 1 \oplus D_0 = 1\)
\item
  New state: \(D_0 \leftarrow 1\), \(D_1 \leftarrow 1\),
  \(D_2 \leftarrow 0\)
\end{itemize}

\begin{Shaded}
\begin{Highlighting}[]
\NormalTok{Input m\_k: 0}
\NormalTok{     |}
\NormalTok{     v}
\NormalTok{0 {-}{-}(+){-}{-}[ D0:1 ]{-}{-}(+){-}{-}[ D1:1 ]{-}{-}(+){-}{-}[ D2:1 ]{-}{-}(+){-}{-}\textgreater{} Output c\_2: 1}
\NormalTok{         (New:1)        (New:1)        (New:0)}
\end{Highlighting}
\end{Shaded}

\textbf{Clock 4:} input \(m_3=1\)

\begin{itemize}
\tightlist
\item
  \(c_3 = 1 \cdot 1 \oplus D_0 = 0\)
\item
  New state: \(D_0 \leftarrow 0\), \(D_1 \leftarrow 0\),
  \(D_2 \leftarrow 1\)
\end{itemize}

\begin{Shaded}
\begin{Highlighting}[]
\NormalTok{Input m\_k: 1}
\NormalTok{     |}
\NormalTok{     v}
\NormalTok{0 {-}{-}(+){-}{-}[ D0:1 ]{-}{-}(+){-}{-}[ D1:1 ]{-}{-}(+){-}{-}[ D2:0 ]{-}{-}(+){-}{-}\textgreater{} Output c\_3: 0}
\NormalTok{         (New:0)        (New:0)        (New:1)}
\end{Highlighting}
\end{Shaded}

\textbf{Clock 5:} flush with input 0

\begin{itemize}
\tightlist
\item
  \(c_4 = 0 \cdot g_0 \oplus D_0 = 0\)
\item
  New state: \(D_0 \leftarrow 0\), \(D_1 \leftarrow 1\),
  \(D_2 \leftarrow 0\)
\end{itemize}

\begin{Shaded}
\begin{Highlighting}[]
\NormalTok{Input m\_k: 0}
\NormalTok{     |}
\NormalTok{     v}
\NormalTok{0 {-}{-}(+){-}{-}[ D0:0 ]{-}{-}(+){-}{-}[ D1:0 ]{-}{-}(+){-}{-}[ D2:1 ]{-}{-}(+){-}{-}\textgreater{} Output c\_4: 0}
\NormalTok{         (New:0)        (New:1)        (New:0)}
\end{Highlighting}
\end{Shaded}

\textbf{Clock 6:} flush again

\begin{itemize}
\tightlist
\item
  \(c_5 = 0 \cdot g_0 \oplus D_0 = 0\)
\item
  New state: \(D_0 \leftarrow 1\), \(D_1 \leftarrow 0\),
  \(D_2 \leftarrow 0\)
\end{itemize}

\begin{Shaded}
\begin{Highlighting}[]
\NormalTok{Input m\_k: 0}
\NormalTok{     |}
\NormalTok{     v}
\NormalTok{0 {-}{-}(+){-}{-}[ D0:0 ]{-}{-}(+){-}{-}[ D1:1 ]{-}{-}(+){-}{-}[ D2:0 ]{-}{-}(+){-}{-}\textgreater{} Output c\_5: 0}
\NormalTok{         (New:1)        (New:0)        (New:0)}
\end{Highlighting}
\end{Shaded}

\textbf{Clock 7:} final output

\begin{itemize}
\tightlist
\item
  \(c_6 = 0 \cdot g_0 \oplus D_0 = 1\)
\item
  All registers return to zero.
\end{itemize}

\begin{Shaded}
\begin{Highlighting}[]
\NormalTok{Input m\_k: 0}
\NormalTok{     |}
\NormalTok{     v}
\NormalTok{0 {-}{-}(+){-}{-}[ D0:1 ]{-}{-}(+){-}{-}[ D1:0 ]{-}{-}(+){-}{-}[ D2:0 ]{-}{-}(+){-}{-}\textgreater{} Output c\_6: 1}
\NormalTok{         (New:0)        (New:0)        (New:0)}
\end{Highlighting}
\end{Shaded}

\textbf{Final result:} \[\mathbf{c}=(1,0,1,0,0,0,1).\] This matches the
polynomial product \[c(X)=(1+X+X^3)(1+X+X^3)=1+X^2+X^6.\]

\subsubsection{Why Does ``Cyclic'' Enable Shift
Registers?}\label{why-does-cyclic-enable-shift-registers}

The key correspondence is
\[\text{cyclic shift} \longleftrightarrow X \cdot c(X) \pmod{X^n-1}.\]

The basic action of a shift register - moving all bits by one position
and injecting feedback - is exactly multiplication by \(X\) followed by
reduction modulo \(X^n-1\). This is why cyclic codes admit LFSR
implementations, whereas general linear codes do not.

In summary:

\begin{Shaded}
\begin{Highlighting}[]
\NormalTok{cyclic{-}shift invariance}
\NormalTok{    {-}\textgreater{} codewords are multiples of g(X)}
\NormalTok{    {-}\textgreater{} encoding becomes polynomial division / multiplication}
\NormalTok{    {-}\textgreater{} implementable by LFSR}
\NormalTok{    {-}\textgreater{} storage O(n{-}k), computation O(n), real{-}time processing}
\end{Highlighting}
\end{Shaded}

This is the fundamental reason cyclic codes became so important, and why
CRC remains the standard error-detection mechanism in protocols such as
Ethernet and USB.


\section{Polynomial Representation of Cyclic Codes}\label{sec:polynomial-representation}

A vector of length \(n\), \[\mathbf{c}=(c_0,c_1,\dots,c_{n-1}),\] can be
mapped to a polynomial of degree at most \(n-1\):

\[c(X)=c_0 + c_1 X + c_2 X^2 + \dots + c_{n-1}X^{n-1}.\]

Under this representation, a one-step cyclic right shift corresponds
exactly to \[X \cdot c(X) \pmod{X^n-1}.\]

Therefore all algebraic operations on a cyclic code naturally live in
the quotient ring \[R_n = \mathbb{F}_q[X]/\langle X^n-1 \rangle.\]

\subsubsection{Example 2.1: Vector \textless-\textgreater{} Polynomial
Conversion}\label{example-2.1-vector---polynomial-conversion}

For \(n=7\):

\begin{longtable}[]{@{}ll@{}}
\toprule\noalign{}
Vector & Polynomial \\
\midrule\noalign{}
\endhead
\bottomrule\noalign{}
\endlastfoot
\((1,0,1,1,0,0,0)\) & \(1+X^2+X^3\) \\
\((0,1,0,0,1,1,0)\) & \(X+X^4+X^5\) \\
\((1,1,1,1,1,1,1)\) & \(1+X+X^2+X^3+X^4+X^5+X^6\) \\
\end{longtable}

\subsubsection{Example 2.2: Cyclic Shift as a Polynomial
Operation}\label{example-2.2-cyclic-shift-as-a-polynomial-operation}

Let \(n=7\) and \(c(X)=1+X^2+X^3\). Then \[Xc(X)=X+X^3+X^4.\] Since the
degree is still below 7, reduction modulo \(X^7-1\) changes nothing. The
corresponding vector is \((0,1,0,1,1,0,0)\), which is exactly the
one-step cyclic shift of \((1,0,1,1,0,0,0)\).

\subsubsection{Example 2.3: When the Highest-Degree Term Wraps
Around}\label{example-2.3-when-the-highest-degree-term-wraps-around}

Let \(n=7\) and \(c(X)=X^4+X^5+X^6\). Then \[Xc(X)=X^5+X^6+X^7.\] Since
\(X^7 \equiv 1 \pmod{X^7-1}\), \[Xc(X) \bmod (X^7-1)=1+X^5+X^6.\] Thus
\[(0,0,0,0,1,1,1) \xrightarrow{\text{1-step shift}} (1,0,0,0,0,1,1).\]


\section{Generator Polynomials}\label{sec:generator-polynomials}

Every cyclic code is completely determined by a unique monic polynomial
\(g(X)\) of \textbf{smallest degree} in the code. This polynomial is
called the \textbf{generator polynomial}.

\textbf{Key properties:}

\begin{itemize}
\tightlist
\item
  \(g(X)\) must divide \(X^n-1\): \[X^n-1 = g(X)h(X).\]
\item
  If \(\deg g = r = n-k\), then the code dimension is \(k\).
\item
  Every codeword is a multiple of \(g(X)\): \[c(X)=m(X)g(X).\]
\end{itemize}

For \textbf{non-systematic encoding}, one simply multiplies the message
polynomial by the generator polynomial: \[c(X)=m(X)g(X).\]

\subsubsection{\texorpdfstring{Example 3.1: Factorization of \(X^7-1\)
and Possible Generator
Polynomials}{Example 3.1: Factorization of X\^{}7-1 and Possible Generator Polynomials}}\label{example-3.1-factorization-of-x7-1-and-possible-generator-polynomials}

Over \(GF(2)\), we have \(-1=1\), so \(X^7-1 = X^7+1\), and it factors
as

\[X^7+1 = (1+X)(1+X+X^3)(1+X^2+X^3).\]

Any divisor of this polynomial can serve as a generator polynomial:

\begin{longtable}[]{@{}
  >{\raggedright\arraybackslash}p{(\columnwidth - 6\tabcolsep) * \real{0.3103}}
  >{\raggedright\arraybackslash}p{(\columnwidth - 6\tabcolsep) * \real{0.1724}}
  >{\raggedright\arraybackslash}p{(\columnwidth - 6\tabcolsep) * \real{0.3103}}
  >{\raggedright\arraybackslash}p{(\columnwidth - 6\tabcolsep) * \real{0.2069}}@{}}
\toprule\noalign{}
\begin{minipage}[b]{\linewidth}\raggedright
\(g(X)\)
\end{minipage} & \begin{minipage}[b]{\linewidth}\raggedright
\(\deg(g)=n-k\)
\end{minipage} & \begin{minipage}[b]{\linewidth}\raggedright
\([n,k]_2\)
\end{minipage} & \begin{minipage}[b]{\linewidth}\raggedright
Number of codewords \(2^k\)
\end{minipage} \\
\midrule\noalign{}
\endhead
\bottomrule\noalign{}
\endlastfoot
\(1 + X\) & 1 & \([7,6]_2\) & 64 \\
\(1 + X + X^3\) & 3 & \([7,4]_2\) & 16 \\
\(1 + X^2 + X^3\) & 3 & \([7,4]_2\) & 16 \\
\((1+X)(1+X+X^3) = 1 + X^2 + X^3 + X^4\) & 4 & \([7,3]_2\) & 8 \\
\((1+X)(1+X^2+X^3) = 1 + X + X^2 + X^4\) & 4 & \([7,3]_2\) & 8 \\
\((1+X+X^3)(1+X^2+X^3) = 1 + X + X^2 + X^3 + X^4 + X^5 + X^6\) & 6 &
\([7,1]_2\) & 2 \\
\(X^7 + 1\) & 7 & \([7,0]_2\) & 1 (trivial) \\
\end{longtable}

\subsubsection{\texorpdfstring{Example 3.2: Non-Systematic Encoding for
a \([7,4]_2\)
Code}{Example 3.2: Non-Systematic Encoding for a {[}7,4{]}\_2 Code}}\label{example-3.2-non-systematic-encoding-for-a-74_2-code}

Let \(g(X)=1+X+X^3\) and let the message be
\[\mathbf{m}=(1,1,0,1), \qquad m(X)=1+X+X^3.\] Then

\[\begin{aligned}
c(X)
&= m(X)g(X) \\
&= (1+X+X^3)(1+X+X^3) \\
&= 1+X+X^3 + X+X^2+X^4 + X^3+X^4+X^6 \\
&= 1 + (1+1)X + X^2 + (1+1)X^3 + (1+1)X^4 + X^6 \\
&= 1 + X^2 + X^6.
\end{aligned}\]

Thus the codeword vector is \[\mathbf{c}=(1,0,1,0,0,0,1).\]

\subsubsection{Example 3.3: Another Non-Systematic Encoding
Example}\label{example-3.3-another-non-systematic-encoding-example}

With the same generator polynomial, let the message be
\[\mathbf{m}=(0,1,0,0), \qquad m(X)=X.\] Then
\[c(X)=X(1+X+X^3)=X+X^2+X^4.\] So the codeword is
\[\mathbf{c}=(0,1,1,0,1,0,0).\] This is exactly the coefficient vector
of \(Xg(X)\), i.e.~the second row of the generator matrix.


\section{Parity-Check Polynomials}\label{sec:parity-check-polynomials}

The quotient obtained by dividing \(X^n-1\) by the generator polynomial
is called the \textbf{parity-check polynomial}:

\[h(X)=\frac{X^n-1}{g(X)}.\]

Its degree is \(k\). To check whether a polynomial \(v(X)\) is a valid
codeword, one tests \[v(X)h(X) \equiv 0 \pmod{X^n-1}.\]

\subsubsection{\texorpdfstring{Example 4.1: Parity-Check Polynomial for
a \([7,4]_2\)
Code}{Example 4.1: Parity-Check Polynomial for a {[}7,4{]}\_2 Code}}\label{example-4.1-parity-check-polynomial-for-a-74_2-code}

If \(g(X)=1+X+X^3\), then \[h(X)=\frac{X^7+1}{1+X+X^3}=1+X+X^2+X^4.\]

Verification: \[g(X)h(X)=(1+X+X^3)(1+X+X^2+X^4)=X^7+1.\] Indeed,
term-by-term cancellation in \(GF(2)\) leaves only \(1+X^7\).

\subsubsection{Example 4.2: Validating a
Codeword}\label{example-4.2-validating-a-codeword}

Consider the codeword \[c(X)=X+X^2+X^4.\] Then
\[c(X)h(X) = (X+X^2+X^4)(1+X+X^2+X^4).\] Expanding and reducing modulo
\(X^7+1\) yields 0, so \(c(X)\) is a valid codeword.


\section{Construction of Generator Matrices}\label{sec:generator-matrices}

Let \[g(X)=g_0+g_1X+\dots+g_{n-k}X^{n-k}.\] Then the generator matrix
\(G\) is built by taking the coefficient vector of \(g(X)\) and shifting
it one position to the right in each row:

\[G = \begin{bmatrix}
g_0 & g_1 & g_2 & \dots & g_{n-k} & 0 & \dots & 0 \\
0 & g_0 & g_1 & \dots & g_{n-k-1} & g_{n-k} & \dots & 0 \\
\vdots & & \ddots & & & & \ddots & \vdots \\
0 & \dots & 0 & g_0 & g_1 & \dots & \dots & g_{n-k}
\end{bmatrix}.\]

The \(i\)-th row corresponds to the coefficient vector of \(X^i g(X)\).

\begin{quote}
\textbf{Why is this a valid generator matrix?} 1. \(g(X)\) is a
codeword, and by cyclic-shift invariance so are
\(Xg(X),X^2g(X),\dots,X^{k-1}g(X)\). 2. These vectors are linearly
independent because their highest nonzero terms occur in different
positions. 3. A \(k\)-dimensional code containing \(k\) linearly
independent codewords has found a basis.
\end{quote}

\subsubsection{\texorpdfstring{Example 5.1: Generator Matrix for a
\([7,4]_2\) Cyclic
Code}{Example 5.1: Generator Matrix for a {[}7,4{]}\_2 Cyclic Code}}\label{example-5.1-generator-matrix-for-a-74_2-cyclic-code}

Let \(g(X)=1+X+X^3\), whose coefficient vector is \((1,1,0,1)\). Then

\[G = \begin{bmatrix}
1 & 1 & 0 & 1 & 0 & 0 & 0 \\
0 & 1 & 1 & 0 & 1 & 0 & 0 \\
0 & 0 & 1 & 1 & 0 & 1 & 0 \\
0 & 0 & 0 & 1 & 1 & 0 & 1
\end{bmatrix}.\]

Interpretation of the rows:

\begin{itemize}
\tightlist
\item
  Row 0: \(g(X)=1+X+X^3\)
\item
  Row 1: \(Xg(X)=X+X^2+X^4\)
\item
  Row 2: \(X^2g(X)=X^2+X^3+X^5\)
\item
  Row 3: \(X^3g(X)=X^3+X^4+X^6\)
\end{itemize}

\subsubsection{Example 5.2: Encoding with the Generator
Matrix}\label{example-5.2-encoding-with-the-generator-matrix}

Let \(\mathbf{m}=(1,0,1,1)\). Then

\[\mathbf{c} = \mathbf{m}G = (1,0,1,1)
\begin{bmatrix}
1 & 1 & 0 & 1 & 0 & 0 & 0 \\
0 & 1 & 1 & 0 & 1 & 0 & 0 \\
0 & 0 & 1 & 1 & 0 & 1 & 0 \\
0 & 0 & 0 & 1 & 1 & 0 & 1
\end{bmatrix}.\]

Carrying out the multiplication over \(GF(2)\) gives
\[\mathbf{c}=(1,1,1,1,1,1,1).\] Thus \[c(X)=1+X+X^2+X^3+X^4+X^5+X^6,\]
and this is divisible by \(g(X)\).

\subsubsection{\texorpdfstring{Example 5.3: A \((7,3)\) Cyclic
Code}{Example 5.3: A (7,3) Cyclic Code}}\label{example-5.3-a-73-cyclic-code}

Let \(g(X)=1+X^2+X^3+X^4\). Then \(k=3\) and

\[G = \begin{bmatrix}
1 & 0 & 1 & 1 & 1 & 0 & 0 \\
0 & 1 & 0 & 1 & 1 & 1 & 0 \\
0 & 0 & 1 & 0 & 1 & 1 & 1
\end{bmatrix}.\]

For message \(\mathbf{m}=(1,1,0)\),
\[\mathbf{c}=(1,1,0)G=(1,1,1,0,0,1,0).\]


\section{Construction of Parity-Check Matrices}\label{sec:parity-check-matrices}

Let \[h(X)=h_0+h_1X+\dots+h_kX^k.\] To construct the parity-check matrix
\(H\), reverse the coefficient order of \(h(X)\) and shift that reversed
vector across the rows:

\[H = \begin{bmatrix}
h_k & h_{k-1} & \dots & h_0 & 0 & \dots & 0 \\
0 & h_k & h_{k-1} & \dots & h_0 & \dots & 0 \\
\vdots & & \ddots & & & \ddots & \vdots \\
0 & \dots & 0 & h_k & h_{k-1} & \dots & h_0
\end{bmatrix}.\]

This matrix satisfies the orthogonality relation \[GH^T=0.\]

\begin{quote}
\textbf{Why reverse the coefficients?}

Since \(g(X)h(X)=X^n-1 \equiv 0 \pmod{X^n-1}\) and every codeword has
the form \(c(X)=m(X)g(X)\), we have \[c(X)h(X) \equiv 0.\] Writing this
condition as inner products of coefficient vectors leads precisely to
the reversed-coefficient construction of \(H\).
\end{quote}

\subsubsection{\texorpdfstring{Example 6.1: Parity-Check Matrix for a
\([7,4]_2\) Cyclic
Code}{Example 6.1: Parity-Check Matrix for a {[}7,4{]}\_2 Cyclic Code}}\label{example-6.1-parity-check-matrix-for-a-74_2-cyclic-code}

For \[h(X)=1+X+X^2+X^4,\] its coefficient vector is \((1,1,1,0,1)\) and
the reversed vector is \((1,0,1,1,1)\). Therefore

\[H = \begin{bmatrix}
1 & 0 & 1 & 1 & 1 & 0 & 0 \\
0 & 1 & 0 & 1 & 1 & 1 & 0 \\
0 & 0 & 1 & 0 & 1 & 1 & 1
\end{bmatrix}.\]

\subsubsection{Example 6.2: Verifying
Orthogonality}\label{example-6.2-verifying-orthogonality}

Take the first row of \(G\), \((1,1,0,1,0,0,0)\), and the first row of
\(H\), \((1,0,1,1,1,0,0)\). Their inner product is
\[1\cdot 1 + 1\cdot 0 + 0\cdot 1 + 1\cdot 1 = 1+1=0 \pmod 2.\] Likewise,
the first row of \(G\) is orthogonal to the second and third rows of
\(H\), and similarly for all remaining rows.

\subsubsection{Example 6.3: Error Detection with the Parity-Check
Matrix}\label{example-6.3-error-detection-with-the-parity-check-matrix}

For \[\mathbf{v}=(1,1,1,1,1,1,1),\] we get
\[H\mathbf{v}^T = \mathbf{0},\] so no error is detected.

If instead \[\mathbf{v}'=(0,1,1,1,1,1,1),\] then
\[H\mathbf{v}'^T = \begin{bmatrix}1\\0\\0\end{bmatrix} \neq \mathbf{0},\]
so the error is detected.


\section{Systematic and Non-Systematic Encoding}\label{sec:systematic-vs-nonsystematic}

Both methods produce valid codewords, but they place the original
message differently inside the codeword.

\begin{longtable}[]{@{}
  >{\raggedright\arraybackslash}p{(\columnwidth - 4\tabcolsep) * \real{0.1132}}
  >{\raggedright\arraybackslash}p{(\columnwidth - 4\tabcolsep) * \real{0.4906}}
  >{\raggedright\arraybackslash}p{(\columnwidth - 4\tabcolsep) * \real{0.3962}}@{}}
\toprule\noalign{}
\begin{minipage}[b]{\linewidth}\raggedright
Category
\end{minipage} & \begin{minipage}[b]{\linewidth}\raggedright
Non-systematic
\end{minipage} & \begin{minipage}[b]{\linewidth}\raggedright
Systematic
\end{minipage} \\
\midrule\noalign{}
\endhead
\bottomrule\noalign{}
\endlastfoot
Encoding rule & \(c(X)=m(X)g(X)\) & \(c(X)=p(X)+X^{n-k}m(X)\) \\
Message position & Mixed throughout the codeword & Preserved explicitly
in fixed positions \\
Ease of decoding & Requires division & Message bits can be read
directly \\
Typical practical use & Rare & Common in real systems \\
\end{longtable}

\subsubsection{Example 7.1: Two Encodings of the Same
Message}\label{example-7.1-two-encodings-of-the-same-message}

Let \(g(X)=1+X+X^3\) and let the message be \(\mathbf{m}=(1,0,1,1)\).

\textbf{Non-systematic encoding:}
\[c(X)=m(X)g(X)=(1+X^2+X^3)(1+X+X^3)=1+X+X^2+X^3+X^4+X^5+X^6.\] The
codeword is \((1,1,1,1,1,1,1)\), and the message is not visibly
preserved as a contiguous block.

\textbf{Systematic encoding:} the resulting codeword has the form
\[(1,0,0,\mathbf{1,0,1,1}),\] so the last four bits are exactly the
original message.


\section{Practical Systematic Encoding}\label{sec:practical-systematic-encoding}

Systematic encoding proceeds in three steps.

\textbf{Step 1. Shift the message polynomial} \[X^{n-k}m(X).\]

\textbf{Step 2. Compute the remainder} \[p(X)=X^{n-k}m(X) \bmod g(X).\]

\textbf{Step 3. Form the codeword} \[c(X)=p(X)+X^{n-k}m(X).\]

Thus the codeword has the structure
\[[\text{parity }(n-k)\text{ bits} \mid \text{message }k\text{ bits}].\]

\begin{quote}
\textbf{Why is this a valid codeword?} Since \(p(X)\) is the remainder
of \(X^{n-k}m(X)\) modulo \(g(X)\), \[X^{n-k}m(X)-p(X)\] is divisible by
\(g(X)\). Over \(GF(2)\), subtraction equals addition, so \(c(X)\) is
also divisible by \(g(X)\).
\end{quote}

\subsubsection{\texorpdfstring{Example 8.1: Systematic Encoding of
\((1,0,1,1)\)}{Example 8.1: Systematic Encoding of (1,0,1,1)}}\label{example-8.1-systematic-encoding-of-1011}

Let \(g(X)=1+X+X^3\), \(n-k=3\), and \(m(X)=1+X^2+X^3\).

\textbf{Step 1:} \[X^3m(X)=X^3+X^5+X^6.\]

\textbf{Step 2:} Divide \(X^6+X^5+X^3\) by \(X^3+X+1\).

\begin{longtable}[]{@{}lll@{}}
\toprule\noalign{}
Step & Dividend & Quotient term \\
\midrule\noalign{}
\endhead
\bottomrule\noalign{}
\endlastfoot
1 & \(X^6 + X^5 + X^3\) & \(X^3\) \\
& \(-(X^6 + X^4 + X^3)\) & \\
& \(= X^5 + X^4\) & \\
2 & \(X^5 + X^4\) & \(X^2\) \\
& \(-(X^5 + X^3 + X^2)\) & \\
& \(= X^4 + X^3 + X^2\) & \\
3 & \(X^4 + X^3 + X^2\) & \(X\) \\
& \(-(X^4 + X^2 + X)\) & \\
& \(= X^3 + X\) & \\
4 & \(X^3 + X\) & \(1\) \\
& \(-(X^3 + X + 1)\) & \\
& \(= 1\) & \\
\end{longtable}

So the quotient is \(X^3+X^2+X+1\), and the remainder is \[p(X)=1.\]

\textbf{Step 3:} \[c(X)=1+X^3+X^5+X^6.\] Thus
\[\mathbf{c}=(\underbrace{1,0,0}_{\text{parity}},\ \underbrace{1,0,1,1}_{\text{message}}).\]

\subsubsection{\texorpdfstring{Example 8.2: Systematic Encoding of
\((1,1,0,0)\)}{Example 8.2: Systematic Encoding of (1,1,0,0)}}\label{example-8.2-systematic-encoding-of-1100}

Here \(m(X)=1+X\).

\textbf{Step 1:} \[X^3m(X)=X^3+X^4.\]

\textbf{Step 2:} Divide by \(X^3+X+1\) to obtain remainder
\[p(X)=X^2+1,\] which corresponds to parity vector \((1,0,1)\).

\textbf{Step 3:} \[c(X)=1+X^2+X^3+X^4.\] So
\[\mathbf{c}=(\underbrace{1,0,1}_{\text{parity}},\ \underbrace{1,1,0,0}_{\text{message}}).\]

\subsubsection{\texorpdfstring{Example 8.3: Systematic Encoding of
\((0,0,0,1)\)}{Example 8.3: Systematic Encoding of (0,0,0,1)}}\label{example-8.3-systematic-encoding-of-0001}

Here \(m(X)=X^3\).

\textbf{Step 1:} \[X^3m(X)=X^6.\]

\textbf{Step 2:} Dividing by \(X^3+X+1\) again yields \[p(X)=X^2+1,\] so
the parity bits are \((1,0,1)\).

\textbf{Step 3:} \[c(X)=1+X^2+X^6.\] Hence
\[\mathbf{c}=(\underbrace{1,0,1}_{\text{parity}},\ \underbrace{0,0,0,1}_{\text{message}}).\]

\subsubsection{\texorpdfstring{Example 8.4: Full Table of Systematic
Codewords for the \([7,4]_2\)
Code}{Example 8.4: Full Table of Systematic Codewords for the {[}7,4{]}\_2 Code}}\label{example-8.4-full-table-of-systematic-codewords-for-the-74_2-code}

\begin{longtable}[]{@{}
  >{\raggedright\arraybackslash}p{(\columnwidth - 6\tabcolsep) * \real{0.2500}}
  >{\raggedright\arraybackslash}p{(\columnwidth - 6\tabcolsep) * \real{0.2500}}
  >{\raggedright\arraybackslash}p{(\columnwidth - 6\tabcolsep) * \real{0.2500}}
  >{\raggedright\arraybackslash}p{(\columnwidth - 6\tabcolsep) * \real{0.2500}}@{}}
\toprule\noalign{}
\begin{minipage}[b]{\linewidth}\raggedright
Message \((m_0m_1m_2m_3)\)
\end{minipage} & \begin{minipage}[b]{\linewidth}\raggedright
\(m(X)\)
\end{minipage} & \begin{minipage}[b]{\linewidth}\raggedright
Parity \((p_0p_1p_2)\)
\end{minipage} & \begin{minipage}[b]{\linewidth}\raggedright
Codeword \((p_0p_1p_2m_0m_1m_2m_3)\)
\end{minipage} \\
\midrule\noalign{}
\endhead
\bottomrule\noalign{}
\endlastfoot
\(0000\) & \(0\) & \(000\) & \(0000000\) \\
\(1000\) & \(1\) & \(110\) & \(1101000\) \\
\(0100\) & \(X\) & \(011\) & \(0110100\) \\
\(1100\) & \(1+X\) & \(101\) & \(1011100\) \\
\(0010\) & \(X^2\) & \(111\) & \(1110010\) \\
\(1010\) & \(1+X^2\) & \(001\) & \(0011010\) \\
\(0110\) & \(X+X^2\) & \(100\) & \(1000110\) \\
\(1110\) & \(1+X+X^2\) & \(010\) & \(0101110\) \\
\(0001\) & \(X^3\) & \(101\) & \(1010001\) \\
\(1001\) & \(1+X^3\) & \(011\) & \(0111001\) \\
\(0101\) & \(X+X^3\) & \(110\) & \(1100101\) \\
\(1101\) & \(1+X+X^3\) & \(000\) & \(0001101\) \\
\(0011\) & \(X^2+X^3\) & \(010\) & \(0100011\) \\
\(1011\) & \(1+X^2+X^3\) & \(100\) & \(1001011\) \\
\(0111\) & \(X+X^2+X^3\) & \(001\) & \(0010111\) \\
\(1111\) & \(1+X+X^2+X^3\) & \(111\) & \(1111111\) \\
\end{longtable}


\section{Syndromes and Error Detection}\label{sec:syndromes-and-error-detection}

The \textbf{remainder} obtained by dividing the received polynomial
\(v(X)\) by the generator polynomial \(g(X)\) is called the
\textbf{syndrome}:

\[s(X)=v(X) \bmod g(X).\]

\begin{itemize}
\tightlist
\item
  If \(s(X)=0\), no error is detected.
\item
  If \(s(X)\neq 0\), an error is detected.
\end{itemize}

\begin{quote}
\textbf{Core principle:} if \(v(X)=c(X)+e(X)\) and \(c(X)\) is a
multiple of \(g(X)\), then \[s(X)=v(X) \bmod g(X)=e(X) \bmod g(X).\]
Thus the syndrome depends only on the error pattern, not on the original
message.
\end{quote}

\subsubsection{Example 9.1: Error-Free
Reception}\label{example-9.1-error-free-reception}

Suppose the codeword \[c(X)=1+X^3+X^5+X^6\] is received without error.
Dividing by \(X^3+X+1\) yields remainder 0.

\begin{longtable}[]{@{}lll@{}}
\toprule\noalign{}
Step & Dividend & Quotient term \\
\midrule\noalign{}
\endhead
\bottomrule\noalign{}
\endlastfoot
1 & \(X^6 + X^5 + X^3 + 1\) & \(X^3\) \\
& \(-(X^6 + X^4 + X^3)\) & \\
& \(= X^5 + X^4 + 1\) & \\
2 & \(X^5 + X^4 + 1\) & \(X^2\) \\
& \(-(X^5 + X^3 + X^2)\) & \\
& \(= X^4 + X^3 + X^2 + 1\) & \\
3 & \(X^4 + X^3 + X^2 + 1\) & \(X\) \\
& \(-(X^4 + X^2 + X)\) & \\
& \(= X^3 + X + 1\) & \\
4 & \(X^3 + X + 1\) & \(1\) \\
& \(-(X^3 + X + 1)\) & \\
& \(= 0\) & \\
\end{longtable}

So the syndrome is 0.

\subsubsection{\texorpdfstring{Example 9.2: A Single-Bit Error at
Position
\(X^2\)}{Example 9.2: A Single-Bit Error at Position X\^{}2}}\label{example-9.2-a-single-bit-error-at-position-x2}

Suppose the received polynomial is \[v(X)=1+X^2+X^3+X^5+X^6,\] so the
error polynomial is \(e(X)=X^2\).

Then \[s(X)=X^2,\] whose vector representation is \((0,0,1)\). Since the
syndrome is nonzero, the error is detected.

\subsubsection{\texorpdfstring{Example 9.3: Complete Syndrome Table for
Single-Bit Errors in the \([7,4]_2\)
Code}{Example 9.3: Complete Syndrome Table for Single-Bit Errors in the {[}7,4{]}\_2 Code}}\label{example-9.3-complete-syndrome-table-for-single-bit-errors-in-the-74_2-code}

For a single-bit error \(e(X)=X^i\), the syndrome is
\(s(X)=X^i \bmod g(X)\).

\begin{longtable}[]{@{}
  >{\raggedright\arraybackslash}p{(\columnwidth - 6\tabcolsep) * \real{0.2500}}
  >{\raggedright\arraybackslash}p{(\columnwidth - 6\tabcolsep) * \real{0.2500}}
  >{\raggedright\arraybackslash}p{(\columnwidth - 6\tabcolsep) * \real{0.2500}}
  >{\raggedright\arraybackslash}p{(\columnwidth - 6\tabcolsep) * \real{0.2500}}@{}}
\toprule\noalign{}
\begin{minipage}[b]{\linewidth}\raggedright
Error position \(i\)
\end{minipage} & \begin{minipage}[b]{\linewidth}\raggedright
\(e(X)\)
\end{minipage} & \begin{minipage}[b]{\linewidth}\raggedright
\(s(X)=e(X) \bmod (1+X+X^3)\)
\end{minipage} & \begin{minipage}[b]{\linewidth}\raggedright
Syndrome vector
\end{minipage} \\
\midrule\noalign{}
\endhead
\bottomrule\noalign{}
\endlastfoot
0 & \(1\) & \(1\) & \((1,0,0)\) \\
1 & \(X\) & \(X\) & \((0,1,0)\) \\
2 & \(X^2\) & \(X^2\) & \((0,0,1)\) \\
3 & \(X^3\) & \(1+X\) & \((1,1,0)\) \\
4 & \(X^4\) & \(X+X^2\) & \((0,1,1)\) \\
5 & \(X^5\) & \(1+X+X^2\) & \((1,1,1)\) \\
6 & \(X^6\) & \(1+X^2\) & \((1,0,1)\) \\
\end{longtable}

Since all 7 syndromes are distinct, the code can correct \textbf{any
single-bit error}.

\subsubsection{Example 9.4: Error Correction Using the
Syndrome}\label{example-9.4-error-correction-using-the-syndrome}

Suppose the received vector is \[\mathbf{v}=(1,0,1,1,0,1,1).\]

\textbf{Step 1.} Compute the syndrome: with \[v(X)=1+X^2+X^3+X^5+X^6,\]
we obtain \[s(X)=X^2,\] so the syndrome vector is \((0,0,1)\).

\textbf{Step 2.} Look up the syndrome table: \((0,0,1)\) corresponds to
error position \(i=2\).

\textbf{Step 3.} Correct the error:
\[\hat{c}(X)=v(X)+X^2 = 1+X^3+X^5+X^6.\] Therefore the corrected
codeword is \[\hat{\mathbf{c}}=(1,0,0,1,0,1,1).\]

\textbf{Step 4.} Extract the message: for a systematic code, the last
four bits are the message, \[\mathbf{m}=(1,0,1,1).\]

\subsubsection{Example 9.5: The Case of a Two-Bit
Error}\label{example-9.5-the-case-of-a-two-bit-error}

Suppose the received vector is \[\mathbf{v}=(0,1,0,1,0,1,1),\] which
corresponds to simultaneous errors at positions 0 and 1. Then
\[e(X)=1+X,\] and the syndrome is \[s(X)=1+X,\] with vector \((1,1,0)\).

But \((1,1,0)\) corresponds in the syndrome table to a \emph{single}
error at position 3. A decoder designed only for single-error correction
will therefore flip the wrong bit, causing a \textbf{miscorrection}.

This is exactly the limitation of the \([7,4]_2\) code: since its
minimum distance is \(d_{\min}=3\), it can correct 1 bit or detect 2
bits, but it cannot reliably correct arbitrary 2-bit errors.


\section{Decoding of Cyclic Codes}\label{sec:decoding}

The details of decoding depend on the specific family of cyclic code
(simple cyclic code, BCH code, RS code, and so on), but the overall
structure is usually the same.

\subsubsection{Step 1: Syndrome
Calculation}\label{step-1-syndrome-calculation}

Given the received polynomial \[v(X)=c(X)+e(X),\] compute
\[s(X)=v(X) \bmod g(X).\] This can be done very efficiently in hardware
by feeding the received stream into an LFSR.

\subsubsection{Step 2: Error-Pattern
Search}\label{step-2-error-pattern-search}

Using the syndrome \(s(X)\), determine the actual error polynomial
\(e(X)\).

\begin{itemize}
\tightlist
\item
  \textbf{Syndrome lookup table:} for short block codes (such as small
  Hamming codes or Golay codes), the syndrome can be mapped directly to
  an error pattern.
\item
  \textbf{Meggitt decoder:} especially hardware-friendly. The syndrome
  register is shifted one clock at a time while logic checks whether the
  current pattern indicates an error in the most significant position.
\item
  \textbf{Algebraic decoding:} for BCH and Reed-Solomon codes, lookup
  tables are impractical. One instead computes the error-locator
  polynomial using methods such as the Berlekamp-Massey algorithm,
  locates roots by Chien search, and recovers magnitudes via the Forney
  algorithm.
\end{itemize}

\subsubsection{Step 3: Error Correction}\label{step-3-error-correction}

Once an estimated error pattern \(\hat{e}(X)\) has been found, subtract
it from the received vector: \[\hat{c}(X)=v(X)-\hat{e}(X).\] For binary
codes, subtraction and addition are the same, so this is just XOR.

\subsubsection{Step 4: Message
Extraction}\label{step-4-message-extraction}

Finally, recover the original message.

\begin{itemize}
\tightlist
\item
  \textbf{Systematic codes:} simply read the designated message
  positions in the codeword.
\item
  \textbf{Non-systematic codes:} divide the corrected codeword by
  \(g(X)\) and take the quotient.
\end{itemize}


\section{Error-Correcting Capability}\label{sec:error-correcting-capability}

A cyclic code is still a linear code, so its error-correcting capability
is determined by its \textbf{minimum distance} \(d_{\min}\). There is no
special correction formula unique to cyclic codes; the same general
relationship for linear codes applies:

\subsubsection{\texorpdfstring{Core Result: \(t\)-Error-Correcting
Capability}{Core Result: t-Error-Correcting Capability}}\label{core-result-t-error-correcting-capability}

A code with minimum distance \(d_{\min}\) can correct up to
\[t = \left\lfloor \frac{d_{\min}-1}{2} \right\rfloor\] errors.

If one is interested only in error \textbf{detection}, then up to
\(d_{\min}-1\) errors can be detected.

\subsubsection{Example 11.1: Error-Correcting Capability of
Representative Cyclic
Codes}\label{example-11.1-error-correcting-capability-of-representative-cyclic-codes}

\begin{longtable}[]{@{}
  >{\raggedright\arraybackslash}p{(\columnwidth - 8\tabcolsep) * \real{0.1132}}
  >{\raggedright\arraybackslash}p{(\columnwidth - 8\tabcolsep) * \real{0.1887}}
  >{\raggedright\arraybackslash}p{(\columnwidth - 8\tabcolsep) * \real{0.2264}}
  >{\raggedright\arraybackslash}p{(\columnwidth - 8\tabcolsep) * \real{0.2642}}
  >{\raggedright\arraybackslash}p{(\columnwidth - 8\tabcolsep) * \real{0.2075}}@{}}
\toprule\noalign{}
\begin{minipage}[b]{\linewidth}\raggedright
Code
\end{minipage} & \begin{minipage}[b]{\linewidth}\raggedright
\([n,k]\)
\end{minipage} & \begin{minipage}[b]{\linewidth}\raggedright
\(d_{\min}\)
\end{minipage} & \begin{minipage}[b]{\linewidth}\raggedright
Correction capability \(t\)
\end{minipage} & \begin{minipage}[b]{\linewidth}\raggedright
Detection capability
\end{minipage} \\
\midrule\noalign{}
\endhead
\bottomrule\noalign{}
\endlastfoot
\([7,4]\) cyclic code (Hamming) & \([7,4]_2\) & \(3\) & \(1\) & \(2\) \\
\([15,7]\) BCH code & \([15,7]\) & \(5\) & \(2\) & \(4\) \\
\([23,12]\) Golay code & \([23,12]\) & \(7\) & \(3\) & \(6\) \\
\([15,5]\) BCH code & \([15,5]\) & \(7\) & \(3\) & \(6\) \\
\([31,21]\) BCH code & \([31,21]\) & \(5\) & \(2\) & \(4\) \\
\end{longtable}

\subsubsection{Limitation of a General Cyclic
Code}\label{limitation-of-a-general-cyclic-code}

For a general cyclic code, there is usually no simple closed-form
formula that gives \(d_{\min}\) directly from \(g(X)\). One must compute
it either by enumerating codewords or by a deeper algebraic analysis.

\subsubsection{BCH Bound: Guaranteed Designed
Distance}\label{bch-bound-guaranteed-designed-distance}

A famous theorem does exist for the important subclass of \textbf{BCH
codes}.

\begin{quote}
\textbf{BCH bound:} Let \(\alpha\) be a primitive \(n\)th root of unity
satisfying \(X^n-1=0\) in \(GF(q^m)\). If the generator polynomial
\(g(X)\) has the consecutive roots
\[\alpha^b,\ \alpha^{b+1},\ \dots,\ \alpha^{b+\delta-2},\] then the
minimum distance satisfies \[d_{\min} \ge \delta.\]
\end{quote}

The parameter \(\delta\) is called the \textbf{designed distance}.
Because of this theorem, one can choose a target correction capability
\(t\) in advance by setting \[\delta = 2t+1.\]

\subsubsection{Example 11.2: Applying the BCH
Bound}\label{example-11.2-applying-the-bch-bound}

For the binary \((15,7)\) BCH code, let \(\alpha\) be a primitive 15th
root of unity in \(GF(2^4)\). If the generator polynomial has roots
\[\alpha, \alpha^2, \alpha^3, \alpha^4,\] then there are 4 consecutive
roots, so \[\delta = 4+1 = 5.\] Therefore
\[d_{\min} \ge 5 \quad \Rightarrow \quad t \ge 2.\] In fact, this code
attains \(d_{\min}=5\), so the BCH bound is tight here.

\subsubsection{Example 11.3: General Cyclic Code vs.~BCH
Code}\label{example-11.3-general-cyclic-code-vs.-bch-code}

For the \([7,4]_2\) cyclic code with \[g(X)=1+X+X^3,\] let \(\alpha\) be
a primitive root of \(X^7+1\) in \(GF(2^3)\). Since
\[g(\alpha)=0, \qquad g(\alpha^2)=0,\] there are two consecutive roots.
Hence the BCH bound gives
\[\delta=3 \quad \Rightarrow \quad d_{\min} \ge 3.\] The actual minimum
distance is exactly 3. In fact this code is both the \([7,4]_2\) Hamming
code and a BCH code.

This ability to design error-correcting performance at the construction
stage is one of the main reasons BCH and Reed-Solomon codes are favored
in practice over arbitrary cyclic codes.


\section{Algebraic Structure of Cyclic Codes}\label{sec:algebraic-structure}

A cyclic code is a \textbf{principal ideal} of the quotient ring
\[R_n = \mathbb{F}_q[X]/\langle X^n-1 \rangle.\]

This means two things:

\begin{enumerate}
\def\labelenumi{\arabic{enumi}.}
\tightlist
\item
  The cyclic code \(\mathcal{C}\) is an ideal of the ring \(R_n\).
\item
  This ideal is generated by a single element \(g(X)\):
  \[\mathcal{C}=\langle g(X) \rangle.\]
\end{enumerate}

\subsubsection{Proof Sketch}\label{proof-sketch}

\textbf{Step 1: Show that \(\mathcal{C}\) is an ideal.}

\begin{itemize}
\tightlist
\item
  \textbf{Closed under addition:} by linearity, if
  \(a(X),b(X)\in\mathcal{C}\) then \(a(X)+b(X)\in\mathcal{C}\).
\item
  \textbf{Absorption:} for any \(r(X)\in R_n\) and
  \(c(X)\in\mathcal{C}\), \[r(X)c(X)\in \mathcal{C}.\] This follows
  because cyclic-shift invariance implies \(Xc(X)\in\mathcal{C}\),
  repeated shifts give \(X^2c(X),X^3c(X),\dots\), and linearity then
  gives closure under any polynomial multiple.
\end{itemize}

\textbf{Step 2: Show that the ideal is principal.}

\begin{itemize}
\tightlist
\item
  Let \(g(X)\) be the monic polynomial of smallest degree in
  \(\mathcal{C}\).
\item
  For any \(c(X)\in\mathcal{C}\), divide:
  \[c(X)=q(X)g(X)+r(X), \qquad \deg r < \deg g.\]
\item
  Since \(r(X)=c(X)-q(X)g(X)\in\mathcal{C}\) and \(g(X)\) has minimal
  degree, we must have \(r(X)=0\).
\item
  Hence every codeword is a multiple of \(g(X)\), so
  \[\mathcal{C}=\langle g(X) \rangle.\]
\end{itemize}

\subsubsection{Example 10.1: Checking the Ideal
Structure}\label{example-10.1-checking-the-ideal-structure}

For the \([7,4]_2\) code with \(g(X)=1+X+X^3\):

\begin{itemize}
\tightlist
\item
  \(g(X)\in\mathcal{C}\)
\item
  \(Xg(X)=X+X^2+X^4\in\mathcal{C}\)
\item
  \((1+X)g(X)=1+X^2+X^3+X^4\in\mathcal{C}\)
\end{itemize}

This illustrates both cyclic closure and absorption.

\subsubsection{\texorpdfstring{Example 10.2: Why Must \(g(X)\) Divide
\(X^n-1\)?}{Example 10.2: Why Must g(X) Divide X\^{}n-1?}}\label{example-10.2-why-must-gx-divide-xn-1}

For \(g(X)=1+X+X^3\), \[X^7+1 = (1+X+X^3)(1+X+X^2+X^4)=g(X)h(X).\] Thus
\(g(X)\) divides \(X^7+1\), which is a necessary condition for any

\section{Ideals in the Quotient Ring and Determination of the Generator Polynomial}\label{sec:ideals-and-generator-polynomials}

From the previous section, cyclic codes correspond one-to-one with
principal ideals of the quotient ring
\[R_n = \mathbb{F}_q[X]/\langle X^n-1 \rangle.\] This naturally leads to
two questions:

\begin{enumerate}
\def\labelenumi{\arabic{enumi}.}
\tightlist
\item
  How many ideals are there?
\item
  How is the generator polynomial that defines each ideal determined?
\end{enumerate}

\subsubsection{Main Theorem: Counting
Ideals}\label{main-theorem-counting-ideals}

Every ideal of \(R_n\) is a principal ideal generated by a monic divisor
\(g(X)\) of \(X^n-1\): \[I = \langle g(X) \rangle.\]

\begin{quote}
\textbf{Proof.}

Consider the canonical projection
\[\pi: \mathbb{F}_q[X] \to R_n, \qquad f(X) \mapsto f(X)+\langle X^n-1 \rangle.\]
This is a surjective ring homomorphism with kernel
\(\langle X^n-1 \rangle\).

By the correspondence theorem for rings, there is a one-to-one
correspondence between ideals \(I\) of \(R_n\) and ideals \(J\) of
\(\mathbb{F}_q[X]\) satisfying \[\langle X^n-1 \rangle \subseteq J.\]
Specifically, \[I = J / \langle X^n-1 \rangle, \qquad J = \pi^{-1}(I).\]

Since \(\mathbb{F}_q[X]\) is a PID (indeed a Euclidean domain), every
ideal \(J\) is principal: \[J=\langle g(X) \rangle\] for a unique monic
polynomial \(g(X)\).

The inclusion \(\langle X^n-1 \rangle \subseteq \langle g(X) \rangle\)
is equivalent to \[g(X) \mid (X^n-1).\]

Conversely, every monic divisor \(g(X)\) of \(X^n-1\) generates an ideal
of \(R_n\).

Therefore there is a one-to-one correspondence:
\[g(X) \mid (X^n-1) \quad \longleftrightarrow \quad \langle g(X) \rangle \text{ is an ideal of } R_n.\]
\(\blacksquare\)
\end{quote}

Hence:

\begin{quote}
\textbf{Number of ideals of \(R_n\) = number of monic divisors of
\(X^n-1\) over \(\mathbb{F}_q[X]\).}
\end{quote}

If \(\gcd(n,q)=1\), so that \(X^n-1\) has no repeated roots, and if
\[X^n-1=f_1(X)f_2(X)\cdots f_s(X)\] is a product of \(s\) distinct
irreducible polynomials, then each factor may be independently included
or excluded from a divisor. Therefore the total number of monic divisors
is \[2^s.\] This count includes both the whole ring (\(g(X)=1\)) and the
zero ideal (\(g(X)=X^n-1\)).

\begin{quote}
\textbf{Why \(2^s\)?} Every monic divisor corresponds uniquely to a
subset \(S \subseteq \{1,2,\dots,s\}\) through
\[g_S(X)=\prod_{i\in S} f_i(X).\] Since the number of subsets is
\(2^s\), the number of monic divisors is also \(2^s\).
\end{quote}

\subsubsection{Procedure for Determining Generator
Polynomials}\label{procedure-for-determining-generator-polynomials}

\textbf{Step 1.} Factor \(X^n-1\) completely into irreducible
polynomials over \(\mathbb{F}_q[X]\).

\textbf{Step 2.} Take products over all subsets of those irreducible
factors.

\textbf{Step 3.} For each generator polynomial \(g(X)\), the
corresponding cyclic code has parameters \[(n,\ n-\deg g(X)).\]

\subsubsection{\texorpdfstring{Example 13.1: Listing All Ideals for
\(n=7\) over
\(\mathbb{F}_2\)}{Example 13.1: Listing All Ideals for n=7 over \textbackslash mathbb\{F\}\_2}}\label{example-13.1-listing-all-ideals-for-n7-over-mathbbf_2}

Since \[X^7+1=(1+X)(1+X+X^3)(1+X^2+X^3),\] there are \(s=3\) distinct
irreducible factors, so the number of ideals is \[2^3=8.\]

\begin{longtable}[]{@{}
  >{\raggedright\arraybackslash}p{(\columnwidth - 10\tabcolsep) * \real{0.1667}}
  >{\raggedright\arraybackslash}p{(\columnwidth - 10\tabcolsep) * \real{0.1667}}
  >{\raggedright\arraybackslash}p{(\columnwidth - 10\tabcolsep) * \real{0.1667}}
  >{\raggedright\arraybackslash}p{(\columnwidth - 10\tabcolsep) * \real{0.1667}}
  >{\raggedright\arraybackslash}p{(\columnwidth - 10\tabcolsep) * \real{0.1667}}
  >{\raggedright\arraybackslash}p{(\columnwidth - 10\tabcolsep) * \real{0.1667}}@{}}
\toprule\noalign{}
\begin{minipage}[b]{\linewidth}\raggedright
Subset \(S\)
\end{minipage} & \begin{minipage}[b]{\linewidth}\raggedright
\(g(X)=\prod_{i\in S} f_i(X)\)
\end{minipage} & \begin{minipage}[b]{\linewidth}\raggedright
\(\deg(g)\)
\end{minipage} & \begin{minipage}[b]{\linewidth}\raggedright
\([n,k]\)
\end{minipage} & \begin{minipage}[b]{\linewidth}\raggedright
Number of codewords \(2^k\)
\end{minipage} & \begin{minipage}[b]{\linewidth}\raggedright
Remark
\end{minipage} \\
\midrule\noalign{}
\endhead
\bottomrule\noalign{}
\endlastfoot
\(\emptyset\) & \(1\) & 0 & \((7,7)\) & 128 & Whole ring (trivial) \\
\(\{1\}\) & \(1+X\) & 1 & \((7,6)\) & 64 & Simple parity-check code \\
\(\{2\}\) & \(1+X+X^3\) & 3 & \([7,4]_2\) & 16 & Hamming \((7,4)\) \\
\(\{3\}\) & \(1+X^2+X^3\) & 3 & \([7,4]_2\) & 16 & \\
\(\{1,2\}\) & \(1+X^2+X^3+X^4\) & 4 & \((7,3)\) & 8 & \\
\(\{1,3\}\) & \(1+X+X^2+X^4\) & 4 & \((7,3)\) & 8 & \\
\(\{2,3\}\) & \(1+X+X^2+X^3+X^4+X^5+X^6\) & 6 & \((7,1)\) & 2 &
Repetition code \\
\(\{1,2,3\}\) & \(X^7+1\) & 7 & \((7,0)\) & 1 & Zero ideal (trivial) \\
\end{longtable}

These 8 ideals form a lattice under inclusion: if
\[g_1(X) \mid g_2(X),\] then
\[\langle g_2(X) \rangle \subseteq \langle g_1(X) \rangle.\] So larger
generator-polynomial degree means smaller code dimension and a smaller
ideal.

\subsubsection{\texorpdfstring{Example 13.2: Number of Ideals for
\(n=15\) over
\(\mathbb{F}_2\)}{Example 13.2: Number of Ideals for n=15 over \textbackslash mathbb\{F\}\_2}}\label{example-13.2-number-of-ideals-for-n15-over-mathbbf_2}

From Exercise 2,
\[X^{15}+1 = (1+X)(1+X+X^2)(1+X+X^4)(1+X^3+X^4)(1+X+X^2+X^3+X^4).\]
There are \(s=5\) irreducible factors, so the number of ideals is
\[2^5=32.\] Among these 32 ideals, there are already several \((15,11)\)
cyclic codes, depending on which degree-4 factor or combination of
factors is selected.

\subsubsection{Example 13.3: Ideal Lattice and Inclusion
Relations}\label{example-13.3-ideal-lattice-and-inclusion-relations}

For \(n=7\), the inclusion lattice looks like this:

\begin{Shaded}
\begin{Highlighting}[]
\NormalTok{        \textless{}1\textgreater{} = R7  (whole ring, k=7)}
\NormalTok{       / |  \textbackslash{}}
\NormalTok{   \textless{}f1\textgreater{} \textless{}f2\textgreater{} \textless{}f3\textgreater{}   (k=6,4,4)}
\NormalTok{     |\textbackslash{} / \textbackslash{} /|}
\NormalTok{  \textless{}f1f2\textgreater{} \textless{}f1f3\textgreater{} \textless{}f2f3\textgreater{}  (k=3,3,1)}
\NormalTok{       \textbackslash{}  |  /}
\NormalTok{      \textless{}f1f2f3\textgreater{} = \textless{}0\textgreater{}  (zero ideal, k=0)}
\end{Highlighting}
\end{Shaded}

As you move upward, the ideals get larger and contain more codewords. As
you move downward, the ideals get smaller, the code dimension decreases,
and the error-correcting capability generally improves.

\subsubsection{\texorpdfstring{Caution When
\(\gcd(n,q) \neq 1\)}{Caution When \textbackslash gcd(n,q) \textbackslash neq 1}}\label{caution-when-gcdnq-neq-1}

If \(\gcd(n,q) \neq 1\), then \(X^n-1\) has repeated roots over
\(\mathbb{F}_q[X]\). For example, over \(\mathbb{F}_2\),
\[X^6+1=(1+X)^2(1+X+X^2)^2,\] since \(\gcd(6,2)=2\).

In that case, one may independently choose exponents
\[0 \le e_i' \le e_i\] for each repeated factor \(f_i(X)^{e_i}\), and
the number of monic divisors becomes \[\prod_{i=1}^s (e_i+1).\] For the
example above, that is \[(2+1)(2+1)=9,\] which is larger than \(2^s=4\).


\section{Practical Importance of Cyclic Codes}\label{sec:practical-importance}

The main practical strength of cyclic codes is that they admit
\textbf{ultra-fast real-time hardware implementations} using LFSRs.

\subsubsection{Principle of the LFSR
Structure}\label{principle-of-the-lfsr-structure}

Whereas a general linear block code requires large matrix
multiplications, a cyclic code reduces to lightweight polynomial
division implemented by:

\begin{itemize}
\tightlist
\item
  \(r=\deg g(X)\) flip-flops,
\item
  a few XOR gates matching the nonzero coefficients of the polynomial,
\item
  and nothing more.
\end{itemize}

\subsubsection{Advantage of Real-Time
Processing}\label{advantage-of-real-time-processing}

\begin{longtable}[]{@{}
  >{\raggedright\arraybackslash}p{(\columnwidth - 2\tabcolsep) * \real{0.5000}}
  >{\raggedright\arraybackslash}p{(\columnwidth - 2\tabcolsep) * \real{0.5000}}@{}}
\toprule\noalign{}
\begin{minipage}[b]{\linewidth}\raggedright
General linear code
\end{minipage} & \begin{minipage}[b]{\linewidth}\raggedright
Cyclic code (LFSR)
\end{minipage} \\
\midrule\noalign{}
\endhead
\bottomrule\noalign{}
\endlastfoot
Buffer the whole block, then multiply by a matrix & Bit-by-bit real-time
processing \\
Large memory required & Only \(r\) registers required \\
High latency & Near-zero latency \\
Complex parallel arithmetic & Simple serial XOR circuit \\
\end{longtable}

\subsubsection{Real Example: CRC in Network
Transmission}\label{real-example-crc-in-network-transmission}

When checking the integrity of gigabit data streams in USB or Ethernet:

\begin{enumerate}
\def\labelenumi{\arabic{enumi}.}
\tightlist
\item
  \textbf{No large memory is needed:} data can be processed as it
  arrives.
\item
  \textbf{On-the-fly computation:} one bit per clock is fed into the
  LFSR.
\item
  \textbf{Immediate result:} when the last bit passes, the remaining
  register contents are the parity bits or syndrome.
\end{enumerate}

\subsubsection{\texorpdfstring{Example 11.1: LFSR Encoder for
\(g(X)=1+X+X^3\)}{Example 11.1: LFSR Encoder for g(X)=1+X+X\^{}3}}\label{example-11.1-lfsr-encoder-for-gx1xx3}

For systematic encoding with \(g(X)=1+X+X^3\), use three flip-flops
\([R_0,R_1,R_2]\). The feedback is the XOR of the incoming bit and the
most significant register \(R_2\). Since the coefficients of \(X^0\) and
\(X^1\) are 1, the feedback is XORed into \(R_0\) and \(R_1\).

Feed the message bits \((1,0,1,1)\) in order from highest degree to
lowest, i.e.~\(m_3,m_2,m_1,m_0\):

\begin{longtable}[]{@{}
  >{\raggedright\arraybackslash}p{(\columnwidth - 10\tabcolsep) * \real{0.1304}}
  >{\raggedright\arraybackslash}p{(\columnwidth - 10\tabcolsep) * \real{0.2391}}
  >{\raggedright\arraybackslash}p{(\columnwidth - 10\tabcolsep) * \real{0.1739}}
  >{\raggedright\arraybackslash}p{(\columnwidth - 10\tabcolsep) * \real{0.1522}}
  >{\raggedright\arraybackslash}p{(\columnwidth - 10\tabcolsep) * \real{0.1522}}
  >{\raggedright\arraybackslash}p{(\columnwidth - 10\tabcolsep) * \real{0.1522}}@{}}
\toprule\noalign{}
\begin{minipage}[b]{\linewidth}\raggedright
Clock
\end{minipage} & \begin{minipage}[b]{\linewidth}\raggedright
Input bit
\end{minipage} & \begin{minipage}[b]{\linewidth}\raggedright
Feedback \(= \text{input} \oplus R_2\)
\end{minipage} & \begin{minipage}[b]{\linewidth}\raggedright
\(R_0\)
\end{minipage} & \begin{minipage}[b]{\linewidth}\raggedright
\(R_1\)
\end{minipage} & \begin{minipage}[b]{\linewidth}\raggedright
\(R_2\)
\end{minipage} \\
\midrule\noalign{}
\endhead
\bottomrule\noalign{}
\endlastfoot
Initial & - & - & 0 & 0 & 0 \\
1 & \(m_3=1\) & \(1 \oplus 0 = 1\) & 1 & 0 & 0 \\
2 & \(m_2=1\) & \(1 \oplus 0 = 1\) & 1 & 1 & 0 \\
3 & \(m_1=0\) & \(0 \oplus 0 = 0\) & 0 & 1 & 1 \\
4 & \(m_0=1\) & \(1 \oplus 1 = 0\) & 0 & 0 & 1 \\
\end{longtable}

After 4 clocks the register state is \((R_0,R_1,R_2)=(0,0,1)\).

\begin{quote}
\textbf{Verification:} In Example 8.1, the remainder for systematic
encoding of \(m(X)=1+X^2+X^3\) was \(p(X)=1\), i.e.~parity vector
\((1,0,0)\). Depending on whether the circuit is interpreted MSB-first
or LSB-first, the register output order may be reversed. Reading
\((0,0,1)\) backward gives \((1,0,0)\).
\end{quote}

\begin{quote}
\textbf{Remark:} The exact behavior of an LFSR depends on the feedback
topology and the bit-input convention. The essential point is that
polynomial division can be executed in real time using only \(n-k\)
registers and a few XOR gates.
\end{quote}


\section{Representative Families of Cyclic Codes}\label{sec:representative-families}

Important code families with cyclic structure include:

\begin{longtable}[]{@{}
  >{\raggedright\arraybackslash}p{(\columnwidth - 4\tabcolsep) * \real{0.2609}}
  >{\raggedright\arraybackslash}p{(\columnwidth - 4\tabcolsep) * \real{0.2609}}
  >{\raggedright\arraybackslash}p{(\columnwidth - 4\tabcolsep) * \real{0.4783}}@{}}
\toprule\noalign{}
\begin{minipage}[b]{\linewidth}\raggedright
Code
\end{minipage} & \begin{minipage}[b]{\linewidth}\raggedright
Characteristic
\end{minipage} & \begin{minipage}[b]{\linewidth}\raggedright
Main applications
\end{minipage} \\
\midrule\noalign{}
\endhead
\bottomrule\noalign{}
\endlastfoot
\textbf{CRC} (Cyclic Redundancy Check) & Error detection only, many
standard generator polynomials & Ethernet, USB, ZIP, storage media \\
\textbf{BCH codes} & Root-based construction guaranteeing
error-correction capability & Flash memory, satellite communication \\
\textbf{Reed-Solomon (RS) codes} & Nonbinary cyclic code, symbol-level
error correction & CD/DVD/Blu-ray, QR codes, deep-space communication \\
\textbf{Some Hamming codes} & Cyclic form for certain parameters & ECC
memory \\
\end{longtable}

\subsubsection{Example 15.1: CRC-8 Generator
Polynomial}\label{example-15.1-crc-8-generator-polynomial}

A standard CRC-8 polynomial is \[g(X)=X^8+X^2+X+1.\]

\begin{itemize}
\tightlist
\item
  Code length: variable (depends on the message length)
\item
  Parity length: 8 bits
\item
  Typical use: sensor networks and embedded systems
\end{itemize}

\subsubsection{Example 15.2: CRC-32 Used in
Ethernet}\label{example-15.2-crc-32-used-in-ethernet}

\[g(X)=X^{32} + X^{26} + X^{23} + X^{22} + X^{16} + X^{12} + X^{11} + X^{10} + X^8 + X^7 + X^5 + X^4 + X^2 + X + 1.\]

\begin{itemize}
\tightlist
\item
  32 parity bits
\item
  Protects frames up to length \(2^{32}-1\) in theory
\item
  Used for the FCS (Frame Check Sequence) in Ethernet
\end{itemize}


\section{Quasi-Cyclic Codes and the Post-Quantum Cryptosystem HQC}\label{sec:qc-and-hqc}

The algebraic structure of cyclic codes plays an important role in
modern cryptography, especially in \textbf{post-quantum cryptography
(PQC)}.

\subsubsection{Quasi-Cyclic Codes (QC)}\label{quasi-cyclic-codes-qc}

If a cyclic code is closed under a 1-step shift, then a
\textbf{quasi-cyclic code} is closed under an \(\ell\)-step shift.

\begin{itemize}
\tightlist
\item
  If the code length is \(n=m\ell\), then shifting any codeword by
  \(\ell\) positions preserves membership in the code.
\item
  If \(\ell=1\), we recover an ordinary cyclic code.
\item
  Generator and parity-check matrices are built from \(m \times m\)
  \textbf{circulant blocks}.
\end{itemize}

\subsubsection{Use in HQC (Hamming
Quasi-Cyclic)}\label{use-in-hqc-hamming-quasi-cyclic}

HQC is a code-based key encapsulation mechanism designed to remain
secure against quantum computers, including attacks based on Shor's
algorithm. It uses quasi-cyclic structure in a central way.

\textbf{The classical problem:} the McEliece cryptosystem is highly
secure but suffers from very large public keys, typically around the
megabyte scale.

\textbf{The HQC idea:} adopt a quasi-cyclic structure in order to
obtain:

\begin{enumerate}
\def\labelenumi{\arabic{enumi}.}
\tightlist
\item
  \textbf{Smaller public keys:} a circulant matrix is determined
  entirely by its first row, so one polynomial can represent a full
  matrix block.
\item
  \textbf{Faster arithmetic:} operations are performed in a polynomial
  ring such as \(\mathbb{F}_2[X]/\langle X^n-1 \rangle\), enabling
  efficient algorithms.
\item
  \textbf{Provable security reductions:} the scheme is related to hard
  syndrome-decoding problems believed to remain difficult even for
  quantum computers.
\end{enumerate}

\subsubsection{Example 16.1: Structure of a Circulant
Block}\label{example-16.1-structure-of-a-circulant-block}

Consider the \(4 \times 4\) circulant matrix whose first row is
\((1,0,1,1)\):

\[C = \begin{bmatrix}
1 & 0 & 1 & 1 \\
1 & 1 & 0 & 1 \\
1 & 1 & 1 & 0 \\
0 & 1 & 1 & 1
\end{bmatrix}.\]

Each row is obtained from the previous one by a one-step cyclic right
shift. Therefore storing only the first row is enough to reconstruct the
entire matrix. This is precisely why quasi-cyclic structure can reduce
public-key size so dramatically.


\section{Exercises}\label{sec:exercises}

\subsubsection{Exercise 1. Determining Whether a Code Is
Cyclic}\label{exercise-1.-determining-whether-a-code-is-cyclic}

Decide whether the following code is cyclic. (\(n=5\))

\[\mathcal{C} = \{00000,\ 10110,\ 01011,\ 11101,\ 01101,\ 10011,\ 00111,\ 11010,\ 11001,\ 00101, \dots\}\]

\textbf{Hint:} apply a one-step cyclic shift to each codeword and check
whether the result is still in the set.


\subsubsection{Exercise 2. Finding Generator
Polynomials}\label{exercise-2.-finding-generator-polynomials}

Suppose that over \(GF(2)\),

\[X^{15} + 1 = (1+X)(1+X+X^2)(1+X+X^4)(1+X^3+X^4)(1+X+X^2+X^3+X^4).\]

\begin{enumerate}
\def\labelenumi{(\alph{enumi})}
\item
  List all degree-4 generator polynomials that can define a \((15,11)\)
  cyclic code.
\item
  If \(g(X)=1+X+X^4\), compute the corresponding parity-check polynomial
  \(h(X)\).
\end{enumerate}


\subsubsection{Exercise 3. Practice with Systematic
Encoding}\label{exercise-3.-practice-with-systematic-encoding}

For the \([7,4]_2\) cyclic code with \(g(X)=1+X+X^3\), systematically
encode the following messages:

\begin{enumerate}
\def\labelenumi{(\alph{enumi})}
\item
  \(\mathbf{m}=(0,1,1,0)\)
\item
  \(\mathbf{m}=(1,1,1,1)\)
\item
  Check that your results agree with the codeword table in Example 8.4.
\end{enumerate}


\subsubsection{Exercise 4. Syndrome Computation and Error
Correction}\label{exercise-4.-syndrome-computation-and-error-correction}

For the \([7,4]_2\) cyclic code, compute the syndrome of each received
vector and correct it if it contains a single-bit error:

\begin{enumerate}
\def\labelenumi{(\alph{enumi})}
\item
  \(\mathbf{v}=(1,1,0,1,0,0,0)\)
\item
  \(\mathbf{v}=(0,1,1,0,1,0,1)\)
\item
  \(\mathbf{v}=(1,0,1,1,1,0,0)\)
\end{enumerate}

\textbf{Hint:} use the syndrome table of Example 9.3.


\subsubsection{Exercise 5. Generator and Parity-Check
Matrices}\label{exercise-5.-generator-and-parity-check-matrices}

Consider the \([7,4]_2\) cyclic code with generator polynomial
\[g(X)=1+X^2+X^3.\]

\begin{enumerate}
\def\labelenumi{(\alph{enumi})}
\item
  Find the non-systematic generator matrix \(G\).
\item
  Compute the parity-check polynomial \(h(X)\).
\item
  Construct the parity-check matrix \(H\).
\item
  Verify that \(GH^T=0\).
\end{enumerate}


\subsubsection{Exercise 6. Understanding the Algebraic
Structure}\label{exercise-6.-understanding-the-algebraic-structure}

\begin{enumerate}
\def\labelenumi{(\alph{enumi})}
\tightlist
\item
  Factor \(X^6+1\) over \(GF(2)\).
\end{enumerate}

\begin{quote}
\textbf{Caution:} since \(\gcd(6,2)=2 \neq 1\), the polynomial \(X^6+1\)
has repeated roots over \(GF(2)\), so some standard results of
cyclic-code theory do not apply in their usual form. Confirm from the
factorization that repeated factors occur.
\end{quote}

\begin{enumerate}
\def\labelenumi{(\alph{enumi})}
\setcounter{enumi}{1}
\item
  Based on that factorization, list the divisors of \(X^6+1\) and
  determine the corresponding \([n,k]\) parameters.
\item
  Find the generator polynomial of the code with largest dimension, and
  determine its minimum distance.
\end{enumerate}


\subsubsection{Exercise 7. Comprehensive
Problem}\label{exercise-7.-comprehensive-problem}

Let the generator polynomial of a \((15,7)\) BCH code be
\[g(X)=1+X^4+X^6+X^7+X^8.\]

\begin{enumerate}
\def\labelenumi{(\alph{enumi})}
\item
  Verify that this is a cyclic code by checking that \(g(X)\) divides
  \(X^{15}+1\).
\item
  Systematically encode the message \[\mathbf{m}=(1,0,0,0,0,0,1).\]
\item
  Verify that the syndrome of the encoded codeword is 0.
\end{enumerate}


%\section{SageMath Practice Code}\label{sec:sagemath-practice-code}
%
%Running the following SageMath script allows you to verify the main
%concepts discussed in this document directly.
%
%\begin{Shaded}
%\begin{Highlighting}[]
%\CommentTok{\# ============================================================}
%\CommentTok{\# Cyclic Codes {-} SageMath practice code}
%\CommentTok{\# Verifies the examples discussed in this document.}
%\CommentTok{\# ============================================================}
%
%\NormalTok{F }\OperatorTok{=}\NormalTok{ GF(}\DecValTok{2}\NormalTok{)}
%\NormalTok{R.}\OperatorTok{\textless{}}\NormalTok{X}\OperatorTok{\textgreater{}} \OperatorTok{=}\NormalTok{ PolynomialRing(F)}
%
%\CommentTok{\# ============================================================}
%\CommentTok{\# 1. Factorization of X\^{}7 + 1 (Example 3.1)}
%\CommentTok{\# ============================================================}
%\BuiltInTok{print}\NormalTok{(}\StringTok{"="} \OperatorTok{*} \DecValTok{70}\NormalTok{)}
%\BuiltInTok{print}\NormalTok{(}\StringTok{"1. Factorization of X\^{}7 + 1 over GF(2)"}\NormalTok{)}
%\BuiltInTok{print}\NormalTok{(}\StringTok{"="} \OperatorTok{*} \DecValTok{70}\NormalTok{)}
%
%\NormalTok{f }\OperatorTok{=}\NormalTok{ X}\OperatorTok{\^{}}\DecValTok{7} \OperatorTok{+} \DecValTok{1}
%\BuiltInTok{print}\NormalTok{(}\SpecialStringTok{f"X\^{}7 + 1 = }\SpecialCharTok{\{}\NormalTok{factor(f)}\SpecialCharTok{\}}\SpecialStringTok{"}\NormalTok{)}
%\BuiltInTok{print}\NormalTok{()}
%
%\CommentTok{\# Three irreducible polynomials}
%\NormalTok{p1 }\OperatorTok{=} \DecValTok{1} \OperatorTok{+}\NormalTok{ X}
%\NormalTok{p2 }\OperatorTok{=} \DecValTok{1} \OperatorTok{+}\NormalTok{ X }\OperatorTok{+}\NormalTok{ X}\OperatorTok{\^{}}\DecValTok{3}
%\NormalTok{p3 }\OperatorTok{=} \DecValTok{1} \OperatorTok{+}\NormalTok{ X}\OperatorTok{\^{}}\DecValTok{2} \OperatorTok{+}\NormalTok{ X}\OperatorTok{\^{}}\DecValTok{3}
%
%\BuiltInTok{print}\NormalTok{(}\StringTok{"Check of irreducible factorization:"}\NormalTok{)}
%\BuiltInTok{print}\NormalTok{(}\SpecialStringTok{f"  (1+X)(1+X+X\^{}3)(1+X\^{}2+X\^{}3) = }\SpecialCharTok{\{}\NormalTok{p1 }\OperatorTok{*}\NormalTok{ p2 }\OperatorTok{*}\NormalTok{ p3}\SpecialCharTok{\}}\SpecialStringTok{"}\NormalTok{)}
%\BuiltInTok{print}\NormalTok{(}\SpecialStringTok{f"  Matches X\^{}7 + 1: }\SpecialCharTok{\{}\StringTok{\textquotesingle{}✓\textquotesingle{}} \ControlFlowTok{if}\NormalTok{ p1 }\OperatorTok{*}\NormalTok{ p2 }\OperatorTok{*}\NormalTok{ p3 }\OperatorTok{==}\NormalTok{ f }\ControlFlowTok{else} \StringTok{\textquotesingle{}✗\textquotesingle{}}\SpecialCharTok{\}}\SpecialStringTok{"}\NormalTok{)}
%\BuiltInTok{print}\NormalTok{()}
%
%\CommentTok{\# List all possible generator polynomials}
%\BuiltInTok{print}\NormalTok{(}\StringTok{"Possible generator polynomials g(X) and corresponding [n, k]:"}\NormalTok{)}
%\BuiltInTok{print}\NormalTok{(}\StringTok{"{-}"} \OperatorTok{*} \DecValTok{50}\NormalTok{)}
%\NormalTok{irreducibles }\OperatorTok{=}\NormalTok{ [p1, p2, p3]}
%\ImportTok{from}\NormalTok{ itertools }\ImportTok{import}\NormalTok{ combinations}
%\NormalTok{all\_divisors }\OperatorTok{=}\NormalTok{ [R(}\DecValTok{1}\NormalTok{)]}
%\ControlFlowTok{for}\NormalTok{ r }\KeywordTok{in} \BuiltInTok{range}\NormalTok{(}\DecValTok{1}\NormalTok{, }\DecValTok{4}\NormalTok{):}
%    \ControlFlowTok{for}\NormalTok{ combo }\KeywordTok{in}\NormalTok{ combinations(irreducibles, r):}
%\NormalTok{        prod }\OperatorTok{=}\NormalTok{ R(}\DecValTok{1}\NormalTok{)}
%        \ControlFlowTok{for}\NormalTok{ p }\KeywordTok{in}\NormalTok{ combo:}
%\NormalTok{            prod }\OperatorTok{*=}\NormalTok{ p}
%\NormalTok{        all\_divisors.append(prod)}
%\NormalTok{all\_divisors.append(f)}
%\NormalTok{all\_divisors.sort(key}\OperatorTok{=}\KeywordTok{lambda}\NormalTok{ p: p.degree())}
%
%\ControlFlowTok{for}\NormalTok{ g }\KeywordTok{in}\NormalTok{ all\_divisors:}
%\NormalTok{    deg }\OperatorTok{=}\NormalTok{ g.degree()}
%\NormalTok{    k }\OperatorTok{=} \DecValTok{7} \OperatorTok{{-}}\NormalTok{ deg}
%    \ControlFlowTok{if} \DecValTok{0} \OperatorTok{\textless{}=}\NormalTok{ k }\OperatorTok{\textless{}=} \DecValTok{7}\NormalTok{:}
%        \BuiltInTok{print}\NormalTok{(}\SpecialStringTok{f"  g(X) = }\SpecialCharTok{\{}\NormalTok{g}\SpecialCharTok{\}}\SpecialStringTok{,  deg = }\SpecialCharTok{\{}\NormalTok{deg}\SpecialCharTok{\}}\SpecialStringTok{,  (7, }\SpecialCharTok{\{}\NormalTok{k}\SpecialCharTok{\}}\SpecialStringTok{),  number of codewords = }\SpecialCharTok{\{}\DecValTok{2}\OperatorTok{\^{}}\NormalTok{k}\SpecialCharTok{\}}\SpecialStringTok{"}\NormalTok{)}
%\BuiltInTok{print}\NormalTok{()}
%
%\CommentTok{\# ============================================================}
%\CommentTok{\# 2. Constructing the (7,4) cyclic code {-} g(X) = 1 + X + X\^{}3}
%\CommentTok{\# ============================================================}
%\BuiltInTok{print}\NormalTok{(}\StringTok{"="} \OperatorTok{*} \DecValTok{70}\NormalTok{)}
%\BuiltInTok{print}\NormalTok{(}\StringTok{"2. [7, 4]\_2 cyclic code: g(X) = 1 + X + X\^{}3"}\NormalTok{)}
%\BuiltInTok{print}\NormalTok{(}\StringTok{"="} \OperatorTok{*} \DecValTok{70}\NormalTok{)}
%
%\NormalTok{g }\OperatorTok{=} \DecValTok{1} \OperatorTok{+}\NormalTok{ X }\OperatorTok{+}\NormalTok{ X}\OperatorTok{\^{}}\DecValTok{3}
%\NormalTok{n }\OperatorTok{=} \DecValTok{7}
%\NormalTok{k }\OperatorTok{=}\NormalTok{ n }\OperatorTok{{-}}\NormalTok{ g.degree()}
%
%\NormalTok{C }\OperatorTok{=}\NormalTok{ codes.CyclicCode(generator\_pol}\OperatorTok{=}\NormalTok{g, length}\OperatorTok{=}\NormalTok{n)}
%\BuiltInTok{print}\NormalTok{(}\SpecialStringTok{f"Code parameters: [}\SpecialCharTok{\{}\NormalTok{C}\SpecialCharTok{.}\NormalTok{length()}\SpecialCharTok{\}}\SpecialStringTok{, }\SpecialCharTok{\{}\NormalTok{C}\SpecialCharTok{.}\NormalTok{dimension()}\SpecialCharTok{\}}\SpecialStringTok{, }\SpecialCharTok{\{}\NormalTok{C}\SpecialCharTok{.}\NormalTok{minimum\_distance()}\SpecialCharTok{\}}\SpecialStringTok{]"}\NormalTok{)}
%\BuiltInTok{print}\NormalTok{(}\SpecialStringTok{f"Generator polynomial: g(X) = }\SpecialCharTok{\{}\NormalTok{g}\SpecialCharTok{\}}\SpecialStringTok{"}\NormalTok{)}
%\BuiltInTok{print}\NormalTok{()}
%
%\CommentTok{\# Parity{-}check polynomial (Example 4.1)}
%\NormalTok{h }\OperatorTok{=}\NormalTok{ (X}\OperatorTok{\^{}}\DecValTok{7} \OperatorTok{+} \DecValTok{1}\NormalTok{) }\OperatorTok{//}\NormalTok{ g}
%\BuiltInTok{print}\NormalTok{(}\SpecialStringTok{f"Parity{-}check polynomial: h(X) = }\SpecialCharTok{\{}\NormalTok{h}\SpecialCharTok{\}}\SpecialStringTok{"}\NormalTok{)}
%\BuiltInTok{print}\NormalTok{(}\SpecialStringTok{f"Verification: g(X) * h(X) = }\SpecialCharTok{\{}\NormalTok{g }\OperatorTok{*}\NormalTok{ h}\SpecialCharTok{\}}\SpecialStringTok{  (== X\^{}7+1: }\SpecialCharTok{\{}\StringTok{\textquotesingle{}✓\textquotesingle{}} \ControlFlowTok{if}\NormalTok{ g }\OperatorTok{*}\NormalTok{ h }\OperatorTok{==}\NormalTok{ f }\ControlFlowTok{else} \StringTok{\textquotesingle{}✗\textquotesingle{}}\SpecialCharTok{\}}\SpecialStringTok{)"}\NormalTok{)}
%\BuiltInTok{print}\NormalTok{()}
%
%\CommentTok{\# ============================================================}
%\CommentTok{\# 3. Generator matrix and parity{-}check matrix}
%\CommentTok{\# ============================================================}
%\BuiltInTok{print}\NormalTok{(}\StringTok{"="} \OperatorTok{*} \DecValTok{70}\NormalTok{)}
%\BuiltInTok{print}\NormalTok{(}\StringTok{"3. Generator matrix G and parity{-}check matrix H"}\NormalTok{)}
%\BuiltInTok{print}\NormalTok{(}\StringTok{"="} \OperatorTok{*} \DecValTok{70}\NormalTok{)}
%
%\NormalTok{G }\OperatorTok{=}\NormalTok{ Matrix(F, k, n)}
%\ControlFlowTok{for}\NormalTok{ i }\KeywordTok{in} \BuiltInTok{range}\NormalTok{(k):}
%\NormalTok{    poly }\OperatorTok{=}\NormalTok{ (X}\OperatorTok{\^{}}\NormalTok{i }\OperatorTok{*}\NormalTok{ g) }\OperatorTok{\%}\NormalTok{ (X}\OperatorTok{\^{}}\NormalTok{n }\OperatorTok{+} \DecValTok{1}\NormalTok{)}
%    \ControlFlowTok{for}\NormalTok{ j }\KeywordTok{in} \BuiltInTok{range}\NormalTok{(n):}
%\NormalTok{        G[i, j] }\OperatorTok{=}\NormalTok{ poly[j]}
%
%\BuiltInTok{print}\NormalTok{(}\StringTok{"Non{-}systematic generator matrix G:"}\NormalTok{)}
%\BuiltInTok{print}\NormalTok{(G)}
%\BuiltInTok{print}\NormalTok{()}
%
%\NormalTok{h\_coeffs }\OperatorTok{=}\NormalTok{ [h[k }\OperatorTok{{-}}\NormalTok{ j] }\ControlFlowTok{for}\NormalTok{ j }\KeywordTok{in} \BuiltInTok{range}\NormalTok{(k }\OperatorTok{+} \DecValTok{1}\NormalTok{)]  }\CommentTok{\# reversed coefficients}
%\NormalTok{H }\OperatorTok{=}\NormalTok{ Matrix(F, n }\OperatorTok{{-}}\NormalTok{ k, n)}
%\ControlFlowTok{for}\NormalTok{ i }\KeywordTok{in} \BuiltInTok{range}\NormalTok{(n }\OperatorTok{{-}}\NormalTok{ k):}
%    \ControlFlowTok{for}\NormalTok{ j }\KeywordTok{in} \BuiltInTok{range}\NormalTok{(k }\OperatorTok{+} \DecValTok{1}\NormalTok{):}
%\NormalTok{        H[i, i }\OperatorTok{+}\NormalTok{ j] }\OperatorTok{=}\NormalTok{ h\_coeffs[j]}
%
%\BuiltInTok{print}\NormalTok{(}\StringTok{"Parity{-}check matrix H:"}\NormalTok{)}
%\BuiltInTok{print}\NormalTok{(H)}
%\BuiltInTok{print}\NormalTok{()}
%
%\NormalTok{product }\OperatorTok{=}\NormalTok{ G }\OperatorTok{*}\NormalTok{ H.transpose()}
%\BuiltInTok{print}\NormalTok{(}\SpecialStringTok{f"Verification that G * H\^{}T = 0: }\SpecialCharTok{\{}\StringTok{\textquotesingle{}✓\textquotesingle{}} \ControlFlowTok{if}\NormalTok{ product }\OperatorTok{==} \DecValTok{0} \ControlFlowTok{else} \StringTok{\textquotesingle{}✗\textquotesingle{}}\SpecialCharTok{\}}\SpecialStringTok{"}\NormalTok{)}
%\BuiltInTok{print}\NormalTok{(G }\OperatorTok{*}\NormalTok{ H.transpose())}
%\BuiltInTok{print}\NormalTok{()}
%
%\CommentTok{\# ============================================================}
%\CommentTok{\# 4. Non{-}systematic encoding (Examples 3.2 and 3.3)}
%\CommentTok{\# ============================================================}
%\BuiltInTok{print}\NormalTok{(}\StringTok{"="} \OperatorTok{*} \DecValTok{70}\NormalTok{)}
%\BuiltInTok{print}\NormalTok{(}\StringTok{"4. Non{-}systematic encoding: c(X) = m(X) * g(X)"}\NormalTok{)}
%\BuiltInTok{print}\NormalTok{(}\StringTok{"="} \OperatorTok{*} \DecValTok{70}\NormalTok{)}
%
%\NormalTok{test\_messages\_nonsys }\OperatorTok{=}\NormalTok{ [}
%\NormalTok{    (}\StringTok{"(1,1,0,1)"}\NormalTok{, }\DecValTok{1} \OperatorTok{+}\NormalTok{ X }\OperatorTok{+}\NormalTok{ X}\OperatorTok{\^{}}\DecValTok{3}\NormalTok{),}
%\NormalTok{    (}\StringTok{"(0,1,0,0)"}\NormalTok{, X),}
%\NormalTok{    (}\StringTok{"(1,0,1,1)"}\NormalTok{, }\DecValTok{1} \OperatorTok{+}\NormalTok{ X}\OperatorTok{\^{}}\DecValTok{2} \OperatorTok{+}\NormalTok{ X}\OperatorTok{\^{}}\DecValTok{3}\NormalTok{),}
%\NormalTok{]}
%
%\ControlFlowTok{for}\NormalTok{ label, m\_poly }\KeywordTok{in}\NormalTok{ test\_messages\_nonsys:}
%\NormalTok{    c\_poly }\OperatorTok{=}\NormalTok{ (m\_poly }\OperatorTok{*}\NormalTok{ g) }\OperatorTok{\%}\NormalTok{ (X}\OperatorTok{\^{}}\NormalTok{n }\OperatorTok{+} \DecValTok{1}\NormalTok{)}
%\NormalTok{    c\_vec }\OperatorTok{=}\NormalTok{ vector(F, [c\_poly[i] }\ControlFlowTok{for}\NormalTok{ i }\KeywordTok{in} \BuiltInTok{range}\NormalTok{(n)])}
%    \BuiltInTok{print}\NormalTok{(}\SpecialStringTok{f"  m = }\SpecialCharTok{\{}\NormalTok{label}\SpecialCharTok{\}}\SpecialStringTok{,  m(X) = }\SpecialCharTok{\{}\NormalTok{m\_poly}\SpecialCharTok{\}}\SpecialStringTok{"}\NormalTok{)}
%    \BuiltInTok{print}\NormalTok{(}\SpecialStringTok{f"  c(X) = }\SpecialCharTok{\{}\NormalTok{c\_poly}\SpecialCharTok{\}}\SpecialStringTok{,  c = }\SpecialCharTok{\{}\BuiltInTok{list}\NormalTok{(c\_vec)}\SpecialCharTok{\}}\SpecialStringTok{"}\NormalTok{)}
%\NormalTok{    remainder }\OperatorTok{=}\NormalTok{ c\_poly }\OperatorTok{\%}\NormalTok{ g}
%    \BuiltInTok{print}\NormalTok{(}\SpecialStringTok{f"  c(X) mod g(X) = }\SpecialCharTok{\{}\NormalTok{remainder}\SpecialCharTok{\}}\SpecialStringTok{  (}\SpecialCharTok{\{}\StringTok{\textquotesingle{}✓ multiple of g(X)\textquotesingle{}} \ControlFlowTok{if}\NormalTok{ remainder }\OperatorTok{==} \DecValTok{0} \ControlFlowTok{else} \StringTok{\textquotesingle{}✗\textquotesingle{}}\SpecialCharTok{\}}\SpecialStringTok{)"}\NormalTok{)}
%    \BuiltInTok{print}\NormalTok{()}
%
%\CommentTok{\# ============================================================}
%\CommentTok{\# 5. Systematic encoding (Examples 8.1{-}8.3)}
%\CommentTok{\# ============================================================}
%\BuiltInTok{print}\NormalTok{(}\StringTok{"="} \OperatorTok{*} \DecValTok{70}\NormalTok{)}
%\BuiltInTok{print}\NormalTok{(}\StringTok{"5. Systematic encoding: c(X) = X\^{}(n{-}k)*m(X) + [X\^{}(n{-}k)*m(X) mod g(X)]"}\NormalTok{)}
%\BuiltInTok{print}\NormalTok{(}\StringTok{"="} \OperatorTok{*} \DecValTok{70}\NormalTok{)}
%
%\KeywordTok{def}\NormalTok{ systematic\_encode(m\_poly, g, n, k):}
%    \CommentTok{"""Perform systematic encoding."""}
%\NormalTok{    r }\OperatorTok{=}\NormalTok{ n }\OperatorTok{{-}}\NormalTok{ k}
%\NormalTok{    shifted }\OperatorTok{=}\NormalTok{ X}\OperatorTok{\^{}}\NormalTok{r }\OperatorTok{*}\NormalTok{ m\_poly}
%\NormalTok{    \_, p }\OperatorTok{=}\NormalTok{ shifted.quo\_rem(g)}
%\NormalTok{    c\_poly }\OperatorTok{=}\NormalTok{ shifted }\OperatorTok{+}\NormalTok{ p}
%\NormalTok{    c\_vec }\OperatorTok{=}\NormalTok{ vector(F, [c\_poly[i] }\ControlFlowTok{for}\NormalTok{ i }\KeywordTok{in} \BuiltInTok{range}\NormalTok{(n)])}
%    \ControlFlowTok{return}\NormalTok{ c\_poly, p, c\_vec}
%
%\NormalTok{test\_messages }\OperatorTok{=}\NormalTok{ [}
%\NormalTok{    (}\StringTok{"(1,0,1,1)"}\NormalTok{, }\DecValTok{1} \OperatorTok{+}\NormalTok{ X}\OperatorTok{\^{}}\DecValTok{2} \OperatorTok{+}\NormalTok{ X}\OperatorTok{\^{}}\DecValTok{3}\NormalTok{),}
%\NormalTok{    (}\StringTok{"(1,1,0,0)"}\NormalTok{, }\DecValTok{1} \OperatorTok{+}\NormalTok{ X),}
%\NormalTok{    (}\StringTok{"(0,0,0,1)"}\NormalTok{, X}\OperatorTok{\^{}}\DecValTok{3}\NormalTok{),}
%\NormalTok{    (}\StringTok{"(0,1,1,0)"}\NormalTok{, X }\OperatorTok{+}\NormalTok{ X}\OperatorTok{\^{}}\DecValTok{2}\NormalTok{),}
%\NormalTok{    (}\StringTok{"(1,1,1,1)"}\NormalTok{, }\DecValTok{1} \OperatorTok{+}\NormalTok{ X }\OperatorTok{+}\NormalTok{ X}\OperatorTok{\^{}}\DecValTok{2} \OperatorTok{+}\NormalTok{ X}\OperatorTok{\^{}}\DecValTok{3}\NormalTok{),}
%\NormalTok{]}
%
%\ControlFlowTok{for}\NormalTok{ label, m\_poly }\KeywordTok{in}\NormalTok{ test\_messages:}
%\NormalTok{    c\_poly, p, c\_vec }\OperatorTok{=}\NormalTok{ systematic\_encode(m\_poly, g, n, k)}
%\NormalTok{    p\_vec }\OperatorTok{=}\NormalTok{ vector(F, [p[i] }\ControlFlowTok{for}\NormalTok{ i }\KeywordTok{in} \BuiltInTok{range}\NormalTok{(n }\OperatorTok{{-}}\NormalTok{ k)])}
%\NormalTok{    m\_vec }\OperatorTok{=}\NormalTok{ vector(F, [m\_poly[i] }\ControlFlowTok{for}\NormalTok{ i }\KeywordTok{in} \BuiltInTok{range}\NormalTok{(k)])}
%    \BuiltInTok{print}\NormalTok{(}\SpecialStringTok{f"  m = }\SpecialCharTok{\{}\NormalTok{label}\SpecialCharTok{\}}\SpecialStringTok{,  m(X) = }\SpecialCharTok{\{}\NormalTok{m\_poly}\SpecialCharTok{\}}\SpecialStringTok{"}\NormalTok{)}
%    \BuiltInTok{print}\NormalTok{(}\SpecialStringTok{f"  parity p(X) = }\SpecialCharTok{\{}\NormalTok{p}\SpecialCharTok{\}}\SpecialStringTok{,  p = }\SpecialCharTok{\{}\BuiltInTok{list}\NormalTok{(p\_vec)}\SpecialCharTok{\}}\SpecialStringTok{"}\NormalTok{)}
%    \BuiltInTok{print}\NormalTok{(}\SpecialStringTok{f"  codeword c = }\SpecialCharTok{\{}\BuiltInTok{list}\NormalTok{(c\_vec)}\SpecialCharTok{\}}\SpecialStringTok{"}\NormalTok{)}
%    \BuiltInTok{print}\NormalTok{(}\SpecialStringTok{f"  structure [parity|message] = }\SpecialCharTok{\{}\BuiltInTok{list}\NormalTok{(p\_vec)}\SpecialCharTok{\}}\SpecialStringTok{|}\SpecialCharTok{\{}\BuiltInTok{list}\NormalTok{(m\_vec)}\SpecialCharTok{\}}\SpecialStringTok{"}\NormalTok{)}
%\NormalTok{    remainder }\OperatorTok{=}\NormalTok{ c\_poly }\OperatorTok{\%}\NormalTok{ g}
%    \BuiltInTok{print}\NormalTok{(}\SpecialStringTok{f"  c(X) mod g(X) = }\SpecialCharTok{\{}\NormalTok{remainder}\SpecialCharTok{\}}\SpecialStringTok{  (}\SpecialCharTok{\{}\StringTok{\textquotesingle{}✓\textquotesingle{}} \ControlFlowTok{if}\NormalTok{ remainder }\OperatorTok{==} \DecValTok{0} \ControlFlowTok{else} \StringTok{\textquotesingle{}✗\textquotesingle{}}\SpecialCharTok{\}}\SpecialStringTok{)"}\NormalTok{)}
%    \BuiltInTok{print}\NormalTok{()}
%
%\CommentTok{\# ============================================================}
%\CommentTok{\# 6. Full systematic codeword table (Example 8.4)}
%\CommentTok{\# ============================================================}
%\BuiltInTok{print}\NormalTok{(}\StringTok{"="} \OperatorTok{*} \DecValTok{70}\NormalTok{)}
%\BuiltInTok{print}\NormalTok{(}\StringTok{"6. Full systematic codeword table for the (7,4) code (16 codewords)"}\NormalTok{)}
%\BuiltInTok{print}\NormalTok{(}\StringTok{"="} \OperatorTok{*} \DecValTok{70}\NormalTok{)}
%\BuiltInTok{print}\NormalTok{(}\SpecialStringTok{f"}\SpecialCharTok{\{}\StringTok{\textquotesingle{}message\textquotesingle{}}\SpecialCharTok{:\textless{}12\}}\SpecialStringTok{ }\SpecialCharTok{\{}\StringTok{\textquotesingle{}m(X)\textquotesingle{}}\SpecialCharTok{:\textless{}20\}}\SpecialStringTok{ }\SpecialCharTok{\{}\StringTok{\textquotesingle{}parity\textquotesingle{}}\SpecialCharTok{:\textless{}10\}}\SpecialStringTok{ }\SpecialCharTok{\{}\StringTok{\textquotesingle{}codeword\textquotesingle{}}\SpecialCharTok{:\textless{}16\}}\SpecialStringTok{ }\SpecialCharTok{\{}\StringTok{\textquotesingle{}wH\textquotesingle{}}\SpecialCharTok{:\textgreater{}4\}}\SpecialStringTok{"}\NormalTok{)}
%\BuiltInTok{print}\NormalTok{(}\StringTok{"{-}"} \OperatorTok{*} \DecValTok{65}\NormalTok{)}
%
%\NormalTok{V4 }\OperatorTok{=}\NormalTok{ VectorSpace(F, }\DecValTok{4}\NormalTok{)}
%\NormalTok{codeword\_set }\OperatorTok{=} \BuiltInTok{set}\NormalTok{()}
%
%\ControlFlowTok{for}\NormalTok{ m\_vec }\KeywordTok{in}\NormalTok{ V4:}
%\NormalTok{    m\_poly }\OperatorTok{=} \BuiltInTok{sum}\NormalTok{(F(m\_vec[i]) }\OperatorTok{*}\NormalTok{ X}\OperatorTok{\^{}}\NormalTok{i }\ControlFlowTok{for}\NormalTok{ i }\KeywordTok{in} \BuiltInTok{range}\NormalTok{(k))}
%\NormalTok{    c\_poly, p, c\_vec }\OperatorTok{=}\NormalTok{ systematic\_encode(m\_poly, g, n, k)}
%\NormalTok{    p\_vec }\OperatorTok{=}\NormalTok{ [p[i] }\ControlFlowTok{for}\NormalTok{ i }\KeywordTok{in} \BuiltInTok{range}\NormalTok{(n }\OperatorTok{{-}}\NormalTok{ k)]}
%\NormalTok{    wH }\OperatorTok{=} \BuiltInTok{sum}\NormalTok{(}\BuiltInTok{int}\NormalTok{(c\_vec[i]) }\ControlFlowTok{for}\NormalTok{ i }\KeywordTok{in} \BuiltInTok{range}\NormalTok{(n))}
%\NormalTok{    codeword\_set.add(}\BuiltInTok{tuple}\NormalTok{(c\_vec))}
%    \BuiltInTok{print}\NormalTok{(}\SpecialStringTok{f"  }\SpecialCharTok{\{}\BuiltInTok{list}\NormalTok{(m\_vec)}\SpecialCharTok{!s:\textless{}12\}}\SpecialStringTok{ }\SpecialCharTok{\{}\BuiltInTok{str}\NormalTok{(m\_poly)}\SpecialCharTok{:\textless{}20\}}\SpecialStringTok{ }\SpecialCharTok{\{}\BuiltInTok{str}\NormalTok{(p\_vec)}\SpecialCharTok{:\textless{}10\}}\SpecialStringTok{ }\SpecialCharTok{\{}\BuiltInTok{list}\NormalTok{(c\_vec)}\SpecialCharTok{!s:\textless{}16\}}\SpecialStringTok{ }\SpecialCharTok{\{}\NormalTok{wH}\SpecialCharTok{:\textgreater{}4\}}\SpecialStringTok{"}\NormalTok{)}
%
%\BuiltInTok{print}\NormalTok{(}\SpecialStringTok{f"}\CharTok{\textbackslash{}n}\SpecialStringTok{Total number of codewords: }\SpecialCharTok{\{}\BuiltInTok{len}\NormalTok{(codeword\_set)}\SpecialCharTok{\}}\SpecialStringTok{ (expected: }\SpecialCharTok{\{}\DecValTok{2}\OperatorTok{\^{}}\NormalTok{k}\SpecialCharTok{\}}\SpecialStringTok{)"}\NormalTok{)}
%\BuiltInTok{print}\NormalTok{()}
%
%\CommentTok{\# ============================================================}
%\CommentTok{\# 7. Verification of cyclic{-}shift invariance}
%\CommentTok{\# ============================================================}
%\BuiltInTok{print}\NormalTok{(}\StringTok{"="} \OperatorTok{*} \DecValTok{70}\NormalTok{)}
%\BuiltInTok{print}\NormalTok{(}\StringTok{"7. Verification of cyclic{-}shift invariance"}\NormalTok{)}
%\BuiltInTok{print}\NormalTok{(}\StringTok{"="} \OperatorTok{*} \DecValTok{70}\NormalTok{)}
%
%\KeywordTok{def}\NormalTok{ cyclic\_shift(v, n):}
%    \CommentTok{"""Shift vector v one step cyclically to the right."""}
%\NormalTok{    v }\OperatorTok{=} \BuiltInTok{list}\NormalTok{(v)}
%    \ControlFlowTok{return} \BuiltInTok{tuple}\NormalTok{([v[}\OperatorTok{{-}}\DecValTok{1}\NormalTok{]] }\OperatorTok{+}\NormalTok{ v[:}\OperatorTok{{-}}\DecValTok{1}\NormalTok{])}
%
%\NormalTok{cyclic\_ok }\OperatorTok{=} \VariableTok{True}
%\ControlFlowTok{for}\NormalTok{ cw }\KeywordTok{in}\NormalTok{ codeword\_set:}
%\NormalTok{    shifted }\OperatorTok{=}\NormalTok{ cw}
%    \ControlFlowTok{for}\NormalTok{ s }\KeywordTok{in} \BuiltInTok{range}\NormalTok{(}\DecValTok{1}\NormalTok{, n):}
%\NormalTok{        shifted }\OperatorTok{=}\NormalTok{ cyclic\_shift(shifted, n)}
%        \ControlFlowTok{if}\NormalTok{ shifted }\KeywordTok{not} \KeywordTok{in}\NormalTok{ codeword\_set:}
%            \BuiltInTok{print}\NormalTok{(}\SpecialStringTok{f"  ✗ }\SpecialCharTok{\{}\NormalTok{cw}\SpecialCharTok{\}}\SpecialStringTok{ {-}\textgreater{} shift by }\SpecialCharTok{\{}\NormalTok{s}\SpecialCharTok{\}}\SpecialStringTok{ {-}\textgreater{} }\SpecialCharTok{\{}\NormalTok{shifted}\SpecialCharTok{\}}\SpecialStringTok{ : not a codeword!"}\NormalTok{)}
%\NormalTok{            cyclic\_ok }\OperatorTok{=} \VariableTok{False}
%
%\ControlFlowTok{if}\NormalTok{ cyclic\_ok:}
%    \BuiltInTok{print}\NormalTok{(}\StringTok{"  ✓ Every cyclic shift of every codeword is still in the code."}\NormalTok{)}
%    \BuiltInTok{print}\NormalTok{(}\StringTok{"  {-}\textgreater{} Confirmed: this is a cyclic code."}\NormalTok{)}
%\BuiltInTok{print}\NormalTok{()}
%
%\CommentTok{\# ============================================================}
%\CommentTok{\# 8. Syndrome table (Example 9.3)}
%\CommentTok{\# ============================================================}
%\BuiltInTok{print}\NormalTok{(}\StringTok{"="} \OperatorTok{*} \DecValTok{70}\NormalTok{)}
%\BuiltInTok{print}\NormalTok{(}\StringTok{"8. Syndrome table for single{-}bit errors"}\NormalTok{)}
%\BuiltInTok{print}\NormalTok{(}\StringTok{"="} \OperatorTok{*} \DecValTok{70}\NormalTok{)}
%\BuiltInTok{print}\NormalTok{(}\SpecialStringTok{f"}\SpecialCharTok{\{}\StringTok{\textquotesingle{}error position i\textquotesingle{}}\SpecialCharTok{:\textless{}16\}}\SpecialStringTok{ }\SpecialCharTok{\{}\StringTok{\textquotesingle{}e(X)\textquotesingle{}}\SpecialCharTok{:\textless{}12\}}\SpecialStringTok{ }\SpecialCharTok{\{}\StringTok{\textquotesingle{}s(X) = e(X) mod g(X)\textquotesingle{}}\SpecialCharTok{:\textless{}24\}}\SpecialStringTok{ }\SpecialCharTok{\{}\StringTok{\textquotesingle{}syndrome vector\textquotesingle{}}\SpecialCharTok{\}}\SpecialStringTok{"}\NormalTok{)}
%\BuiltInTok{print}\NormalTok{(}\StringTok{"{-}"} \OperatorTok{*} \DecValTok{70}\NormalTok{)}
%
%\NormalTok{syndrome\_table }\OperatorTok{=}\NormalTok{ \{\}}
%\ControlFlowTok{for}\NormalTok{ i }\KeywordTok{in} \BuiltInTok{range}\NormalTok{(n):}
%\NormalTok{    e\_poly }\OperatorTok{=}\NormalTok{ X}\OperatorTok{\^{}}\NormalTok{i}
%\NormalTok{    s\_poly }\OperatorTok{=}\NormalTok{ e\_poly }\OperatorTok{\%}\NormalTok{ g}
%\NormalTok{    s\_vec }\OperatorTok{=} \BuiltInTok{tuple}\NormalTok{(F(s\_poly[j]) }\ControlFlowTok{for}\NormalTok{ j }\KeywordTok{in} \BuiltInTok{range}\NormalTok{(n }\OperatorTok{{-}}\NormalTok{ k))}
%\NormalTok{    syndrome\_table[s\_vec] }\OperatorTok{=}\NormalTok{ i}
%    \BuiltInTok{print}\NormalTok{(}\SpecialStringTok{f"  }\SpecialCharTok{\{}\NormalTok{i}\SpecialCharTok{:\textless{}16\}}\SpecialStringTok{ }\SpecialCharTok{\{}\BuiltInTok{str}\NormalTok{(e\_poly)}\SpecialCharTok{:\textless{}12\}}\SpecialStringTok{ }\SpecialCharTok{\{}\BuiltInTok{str}\NormalTok{(s\_poly)}\SpecialCharTok{:\textless{}24\}}\SpecialStringTok{ }\SpecialCharTok{\{}\NormalTok{s\_vec}\SpecialCharTok{\}}\SpecialStringTok{"}\NormalTok{)}
%
%\BuiltInTok{print}\NormalTok{(}\SpecialStringTok{f"}\CharTok{\textbackslash{}n}\SpecialStringTok{Number of distinct syndromes: }\SpecialCharTok{\{}\BuiltInTok{len}\NormalTok{(syndrome\_table)}\SpecialCharTok{\}}\SpecialStringTok{ (expected: }\SpecialCharTok{\{}\NormalTok{n}\SpecialCharTok{\}}\SpecialStringTok{)"}\NormalTok{)}
%\NormalTok{all\_distinct }\OperatorTok{=} \BuiltInTok{len}\NormalTok{(syndrome\_table) }\OperatorTok{==}\NormalTok{ n}
%\BuiltInTok{print}\NormalTok{(}\SpecialStringTok{f"All single{-}bit errors are correctable: }\SpecialCharTok{\{}\StringTok{\textquotesingle{}✓\textquotesingle{}} \ControlFlowTok{if}\NormalTok{ all\_distinct }\ControlFlowTok{else} \StringTok{\textquotesingle{}✗\textquotesingle{}}\SpecialCharTok{\}}\SpecialStringTok{"}\NormalTok{)}
%\BuiltInTok{print}\NormalTok{()}
%
%\CommentTok{\# ============================================================}
%\CommentTok{\# 9. Syndrome{-}based error{-}correction simulation (Example 9.4)}
%\CommentTok{\# ============================================================}
%\BuiltInTok{print}\NormalTok{(}\StringTok{"="} \OperatorTok{*} \DecValTok{70}\NormalTok{)}
%\BuiltInTok{print}\NormalTok{(}\StringTok{"9. Syndrome{-}based error detection and correction simulation"}\NormalTok{)}
%\BuiltInTok{print}\NormalTok{(}\StringTok{"="} \OperatorTok{*} \DecValTok{70}\NormalTok{)}
%
%\NormalTok{m\_orig }\OperatorTok{=} \DecValTok{1} \OperatorTok{+}\NormalTok{ X}\OperatorTok{\^{}}\DecValTok{2} \OperatorTok{+}\NormalTok{ X}\OperatorTok{\^{}}\DecValTok{3}  \CommentTok{\# m = (1,0,1,1)}
%\NormalTok{c\_poly\_orig, \_, c\_vec\_orig }\OperatorTok{=}\NormalTok{ systematic\_encode(m\_orig, g, n, k)}
%\BuiltInTok{print}\NormalTok{(}\SpecialStringTok{f"Original message: m = (1,0,1,1)"}\NormalTok{)}
%\BuiltInTok{print}\NormalTok{(}\SpecialStringTok{f"Codeword:         c = }\SpecialCharTok{\{}\BuiltInTok{list}\NormalTok{(c\_vec\_orig)}\SpecialCharTok{\}}\SpecialStringTok{"}\NormalTok{)}
%\BuiltInTok{print}\NormalTok{()}
%
%\BuiltInTok{print}\NormalTok{(}\StringTok{"Inject a single{-}bit error {-}\textgreater{} compute syndrome {-}\textgreater{} correct:"}\NormalTok{)}
%\BuiltInTok{print}\NormalTok{(}\StringTok{"{-}"} \OperatorTok{*} \DecValTok{70}\NormalTok{)}
%
%\ControlFlowTok{for}\NormalTok{ err\_pos }\KeywordTok{in} \BuiltInTok{range}\NormalTok{(n):}
%\NormalTok{    e\_poly }\OperatorTok{=}\NormalTok{ X}\OperatorTok{\^{}}\NormalTok{err\_pos}
%\NormalTok{    v\_poly }\OperatorTok{=}\NormalTok{ c\_poly\_orig }\OperatorTok{+}\NormalTok{ e\_poly}
%\NormalTok{    v\_vec }\OperatorTok{=}\NormalTok{ vector(F, [v\_poly[i] }\ControlFlowTok{for}\NormalTok{ i }\KeywordTok{in} \BuiltInTok{range}\NormalTok{(n)])}
%
%\NormalTok{    s\_poly }\OperatorTok{=}\NormalTok{ v\_poly }\OperatorTok{\%}\NormalTok{ g}
%\NormalTok{    s\_vec }\OperatorTok{=} \BuiltInTok{tuple}\NormalTok{(F(s\_poly[j]) }\ControlFlowTok{for}\NormalTok{ j }\KeywordTok{in} \BuiltInTok{range}\NormalTok{(n }\OperatorTok{{-}}\NormalTok{ k))}
%
%    \ControlFlowTok{if}\NormalTok{ s\_vec }\KeywordTok{in}\NormalTok{ syndrome\_table:}
%\NormalTok{        est\_pos }\OperatorTok{=}\NormalTok{ syndrome\_table[s\_vec]}
%\NormalTok{        corrected\_poly }\OperatorTok{=}\NormalTok{ v\_poly }\OperatorTok{+}\NormalTok{ X}\OperatorTok{\^{}}\NormalTok{est\_pos}
%\NormalTok{        corrected\_vec }\OperatorTok{=}\NormalTok{ vector(F, [corrected\_poly[i] }\ControlFlowTok{for}\NormalTok{ i }\KeywordTok{in} \BuiltInTok{range}\NormalTok{(n)])}
%\NormalTok{        success }\OperatorTok{=}\NormalTok{ (corrected\_vec }\OperatorTok{==}\NormalTok{ c\_vec\_orig)}
%        \BuiltInTok{print}\NormalTok{(}\SpecialStringTok{f"  error position=}\SpecialCharTok{\{}\NormalTok{err\_pos}\SpecialCharTok{\}}\SpecialStringTok{, received=}\SpecialCharTok{\{}\BuiltInTok{list}\NormalTok{(v\_vec)}\SpecialCharTok{\}}\SpecialStringTok{, "}
%              \SpecialStringTok{f"syndrome=}\SpecialCharTok{\{}\NormalTok{s\_vec}\SpecialCharTok{\}}\SpecialStringTok{, estimated position=}\SpecialCharTok{\{}\NormalTok{est\_pos}\SpecialCharTok{\}}\SpecialStringTok{, "}
%              \SpecialStringTok{f"corrected=}\SpecialCharTok{\{}\StringTok{\textquotesingle{}✓\textquotesingle{}} \ControlFlowTok{if}\NormalTok{ success }\ControlFlowTok{else} \StringTok{\textquotesingle{}✗\textquotesingle{}}\SpecialCharTok{\}}\SpecialStringTok{"}\NormalTok{)}
%    \ControlFlowTok{else}\NormalTok{:}
%        \BuiltInTok{print}\NormalTok{(}\SpecialStringTok{f"  error position=}\SpecialCharTok{\{}\NormalTok{err\_pos}\SpecialCharTok{\}}\SpecialStringTok{, syndrome=}\SpecialCharTok{\{}\NormalTok{s\_vec}\SpecialCharTok{\}}\SpecialStringTok{, not in table!"}\NormalTok{)}
%
%\BuiltInTok{print}\NormalTok{()}
%
%\CommentTok{\# ============================================================}
%\CommentTok{\# 10. Miscorrection under a 2{-}bit error (Example 9.5)}
%\CommentTok{\# ============================================================}
%\BuiltInTok{print}\NormalTok{(}\StringTok{"="} \OperatorTok{*} \DecValTok{70}\NormalTok{)}
%\BuiltInTok{print}\NormalTok{(}\StringTok{"10. Miscorrection under a 2{-}bit error"}\NormalTok{)}
%\BuiltInTok{print}\NormalTok{(}\StringTok{"="} \OperatorTok{*} \DecValTok{70}\NormalTok{)}
%
%\NormalTok{e\_poly\_2bit }\OperatorTok{=} \DecValTok{1} \OperatorTok{+}\NormalTok{ X}
%\NormalTok{v\_poly\_2bit }\OperatorTok{=}\NormalTok{ c\_poly\_orig }\OperatorTok{+}\NormalTok{ e\_poly\_2bit}
%\NormalTok{v\_vec\_2bit }\OperatorTok{=}\NormalTok{ vector(F, [v\_poly\_2bit[i] }\ControlFlowTok{for}\NormalTok{ i }\KeywordTok{in} \BuiltInTok{range}\NormalTok{(n)])}
%\NormalTok{s\_poly\_2bit }\OperatorTok{=}\NormalTok{ v\_poly\_2bit }\OperatorTok{\%}\NormalTok{ g}
%\NormalTok{s\_vec\_2bit }\OperatorTok{=} \BuiltInTok{tuple}\NormalTok{(F(s\_poly\_2bit[j]) }\ControlFlowTok{for}\NormalTok{ j }\KeywordTok{in} \BuiltInTok{range}\NormalTok{(n }\OperatorTok{{-}}\NormalTok{ k))}
%
%\BuiltInTok{print}\NormalTok{(}\SpecialStringTok{f"Original codeword: }\SpecialCharTok{\{}\BuiltInTok{list}\NormalTok{(c\_vec\_orig)}\SpecialCharTok{\}}\SpecialStringTok{"}\NormalTok{)}
%\BuiltInTok{print}\NormalTok{(}\SpecialStringTok{f"Error pattern:     e(X) = 1 + X (positions 0 and 1)"}\NormalTok{)}
%\BuiltInTok{print}\NormalTok{(}\SpecialStringTok{f"Received vector:   }\SpecialCharTok{\{}\BuiltInTok{list}\NormalTok{(v\_vec\_2bit)}\SpecialCharTok{\}}\SpecialStringTok{"}\NormalTok{)}
%\BuiltInTok{print}\NormalTok{(}\SpecialStringTok{f"Syndrome:          }\SpecialCharTok{\{}\NormalTok{s\_vec\_2bit}\SpecialCharTok{\}}\SpecialStringTok{"}\NormalTok{)}
%
%\ControlFlowTok{if}\NormalTok{ s\_vec\_2bit }\KeywordTok{in}\NormalTok{ syndrome\_table:}
%\NormalTok{    est\_pos }\OperatorTok{=}\NormalTok{ syndrome\_table[s\_vec\_2bit]}
%\NormalTok{    corrected }\OperatorTok{=}\NormalTok{ v\_poly\_2bit }\OperatorTok{+}\NormalTok{ X}\OperatorTok{\^{}}\NormalTok{est\_pos}
%\NormalTok{    corrected\_vec }\OperatorTok{=}\NormalTok{ vector(F, [corrected[i] }\ControlFlowTok{for}\NormalTok{ i }\KeywordTok{in} \BuiltInTok{range}\NormalTok{(n)])}
%\NormalTok{    is\_correct }\OperatorTok{=}\NormalTok{ (corrected\_vec }\OperatorTok{==}\NormalTok{ c\_vec\_orig)}
%    \BuiltInTok{print}\NormalTok{(}\SpecialStringTok{f"Estimated error position: }\SpecialCharTok{\{}\NormalTok{est\_pos}\SpecialCharTok{\}}\SpecialStringTok{ (actual positions: 0 and 1)"}\NormalTok{)}
%    \BuiltInTok{print}\NormalTok{(}\SpecialStringTok{f"Correction result: }\SpecialCharTok{\{}\BuiltInTok{list}\NormalTok{(corrected\_vec)}\SpecialCharTok{\}}\SpecialStringTok{"}\NormalTok{)}
%    \BuiltInTok{print}\NormalTok{(}\SpecialStringTok{f"Successful correction: }\SpecialCharTok{\{}\StringTok{\textquotesingle{}✓\textquotesingle{}} \ControlFlowTok{if}\NormalTok{ is\_correct }\ControlFlowTok{else} \StringTok{\textquotesingle{}✗ miscorrection occurred!\textquotesingle{}}\SpecialCharTok{\}}\SpecialStringTok{"}\NormalTok{)}
%\BuiltInTok{print}\NormalTok{()}
%
%\CommentTok{\# ============================================================}
%\CommentTok{\# 11. Codeword validation using the parity{-}check polynomial}
%\CommentTok{\# ============================================================}
%\BuiltInTok{print}\NormalTok{(}\StringTok{"="} \OperatorTok{*} \DecValTok{70}\NormalTok{)}
%\BuiltInTok{print}\NormalTok{(}\StringTok{"11. Validating codewords with h(X): c(X)*h(X) ≡ 0 mod (X\^{}7+1)"}\NormalTok{)}
%\BuiltInTok{print}\NormalTok{(}\StringTok{"="} \OperatorTok{*} \DecValTok{70}\NormalTok{)}
%
%\NormalTok{h }\OperatorTok{=}\NormalTok{ (X}\OperatorTok{\^{}}\DecValTok{7} \OperatorTok{+} \DecValTok{1}\NormalTok{) }\OperatorTok{//}\NormalTok{ g}
%\BuiltInTok{print}\NormalTok{(}\SpecialStringTok{f"h(X) = }\SpecialCharTok{\{}\NormalTok{h}\SpecialCharTok{\}}\SpecialStringTok{"}\NormalTok{)}
%\BuiltInTok{print}\NormalTok{()}
%
%\NormalTok{all\_valid }\OperatorTok{=} \VariableTok{True}
%\ControlFlowTok{for}\NormalTok{ cw }\KeywordTok{in}\NormalTok{ codeword\_set:}
%\NormalTok{    c\_poly\_check }\OperatorTok{=} \BuiltInTok{sum}\NormalTok{(F(cw[i]) }\OperatorTok{*}\NormalTok{ X}\OperatorTok{\^{}}\NormalTok{i }\ControlFlowTok{for}\NormalTok{ i }\KeywordTok{in} \BuiltInTok{range}\NormalTok{(n))}
%\NormalTok{    product }\OperatorTok{=}\NormalTok{ (c\_poly\_check }\OperatorTok{*}\NormalTok{ h) }\OperatorTok{\%}\NormalTok{ (X}\OperatorTok{\^{}}\NormalTok{n }\OperatorTok{+} \DecValTok{1}\NormalTok{)}
%    \ControlFlowTok{if}\NormalTok{ product }\OperatorTok{!=} \DecValTok{0}\NormalTok{:}
%        \BuiltInTok{print}\NormalTok{(}\SpecialStringTok{f"  ✗ c = }\SpecialCharTok{\{}\NormalTok{cw}\SpecialCharTok{\}}\SpecialStringTok{: c(X)*h(X) mod (X\^{}7+1) = }\SpecialCharTok{\{}\NormalTok{product}\SpecialCharTok{\}}\SpecialStringTok{"}\NormalTok{)}
%\NormalTok{        all\_valid }\OperatorTok{=} \VariableTok{False}
%
%\ControlFlowTok{if}\NormalTok{ all\_valid:}
%    \BuiltInTok{print}\NormalTok{(}\SpecialStringTok{f"  ✓ For all }\SpecialCharTok{\{}\BuiltInTok{len}\NormalTok{(codeword\_set)}\SpecialCharTok{\}}\SpecialStringTok{ codewords, c(X)*h(X) ≡ 0 (mod X\^{}7+1)"}\NormalTok{)}
%\BuiltInTok{print}\NormalTok{()}
%
%\CommentTok{\# ============================================================}
%\CommentTok{\# 12. Cyclic codes generated by different divisors of X\^{}7+1}
%\CommentTok{\# ============================================================}
%\BuiltInTok{print}\NormalTok{(}\StringTok{"="} \OperatorTok{*} \DecValTok{70}\NormalTok{)}
%\BuiltInTok{print}\NormalTok{(}\StringTok{"12. Cyclic codes generated by various divisors of X\^{}7+1"}\NormalTok{)}
%\BuiltInTok{print}\NormalTok{(}\StringTok{"="} \OperatorTok{*} \DecValTok{70}\NormalTok{)}
%
%\NormalTok{generators }\OperatorTok{=}\NormalTok{ [}
%\NormalTok{    (}\StringTok{"1+X"}\NormalTok{,                    }\DecValTok{1} \OperatorTok{+}\NormalTok{ X),}
%\NormalTok{    (}\StringTok{"1+X+X\^{}3"}\NormalTok{,                }\DecValTok{1} \OperatorTok{+}\NormalTok{ X }\OperatorTok{+}\NormalTok{ X}\OperatorTok{\^{}}\DecValTok{3}\NormalTok{),}
%\NormalTok{    (}\StringTok{"1+X\^{}2+X\^{}3"}\NormalTok{,              }\DecValTok{1} \OperatorTok{+}\NormalTok{ X}\OperatorTok{\^{}}\DecValTok{2} \OperatorTok{+}\NormalTok{ X}\OperatorTok{\^{}}\DecValTok{3}\NormalTok{),}
%\NormalTok{    (}\StringTok{"(1+X)(1+X+X\^{}3)"}\NormalTok{,         (}\DecValTok{1} \OperatorTok{+}\NormalTok{ X) }\OperatorTok{*}\NormalTok{ (}\DecValTok{1} \OperatorTok{+}\NormalTok{ X }\OperatorTok{+}\NormalTok{ X}\OperatorTok{\^{}}\DecValTok{3}\NormalTok{)),}
%\NormalTok{    (}\StringTok{"(1+X)(1+X\^{}2+X\^{}3)"}\NormalTok{,       (}\DecValTok{1} \OperatorTok{+}\NormalTok{ X) }\OperatorTok{*}\NormalTok{ (}\DecValTok{1} \OperatorTok{+}\NormalTok{ X}\OperatorTok{\^{}}\DecValTok{2} \OperatorTok{+}\NormalTok{ X}\OperatorTok{\^{}}\DecValTok{3}\NormalTok{)),}
%\NormalTok{    (}\StringTok{"(1+X+X\^{}3)(1+X\^{}2+X\^{}3)"}\NormalTok{,   (}\DecValTok{1} \OperatorTok{+}\NormalTok{ X }\OperatorTok{+}\NormalTok{ X}\OperatorTok{\^{}}\DecValTok{3}\NormalTok{) }\OperatorTok{*}\NormalTok{ (}\DecValTok{1} \OperatorTok{+}\NormalTok{ X}\OperatorTok{\^{}}\DecValTok{2} \OperatorTok{+}\NormalTok{ X}\OperatorTok{\^{}}\DecValTok{3}\NormalTok{)),}
%\NormalTok{]}
%
%\BuiltInTok{print}\NormalTok{(}\SpecialStringTok{f"}\SpecialCharTok{\{}\StringTok{\textquotesingle{}g(X)\textquotesingle{}}\SpecialCharTok{:\textless{}30\}}\SpecialStringTok{ }\SpecialCharTok{\{}\StringTok{\textquotesingle{}deg\textquotesingle{}}\SpecialCharTok{:\textless{}6\}}\SpecialStringTok{ }\SpecialCharTok{\{}\StringTok{\textquotesingle{}(n,k)\textquotesingle{}}\SpecialCharTok{:\textless{}10\}}\SpecialStringTok{ }\SpecialCharTok{\{}\StringTok{\textquotesingle{}d\_min\textquotesingle{}}\SpecialCharTok{:\textless{}8\}}\SpecialStringTok{ }\SpecialCharTok{\{}\StringTok{\textquotesingle{}\# codewords\textquotesingle{}}\SpecialCharTok{\}}\SpecialStringTok{"}\NormalTok{)}
%\BuiltInTok{print}\NormalTok{(}\StringTok{"{-}"} \OperatorTok{*} \DecValTok{70}\NormalTok{)}
%
%\ControlFlowTok{for}\NormalTok{ label, gen }\KeywordTok{in}\NormalTok{ generators:}
%    \ControlFlowTok{try}\NormalTok{:}
%\NormalTok{        C\_test }\OperatorTok{=}\NormalTok{ codes.CyclicCode(generator\_pol}\OperatorTok{=}\NormalTok{gen, length}\OperatorTok{=}\DecValTok{7}\NormalTok{)}
%\NormalTok{        d\_min }\OperatorTok{=}\NormalTok{ C\_test.minimum\_distance()}
%\NormalTok{        dim }\OperatorTok{=}\NormalTok{ C\_test.dimension()}
%        \BuiltInTok{print}\NormalTok{(}\SpecialStringTok{f"  }\SpecialCharTok{\{}\NormalTok{label}\SpecialCharTok{:\textless{}30\}}\SpecialStringTok{ }\SpecialCharTok{\{}\NormalTok{gen}\SpecialCharTok{.}\NormalTok{degree()}\SpecialCharTok{:\textless{}6\}}\SpecialStringTok{ (7,}\SpecialCharTok{\{}\NormalTok{dim}\SpecialCharTok{\}}\SpecialStringTok{)}\SpecialCharTok{\{}\StringTok{\textquotesingle{}\textquotesingle{}}\SpecialCharTok{:\textless{}5\}}\SpecialStringTok{ }\SpecialCharTok{\{}\NormalTok{d\_min}\SpecialCharTok{:\textless{}8\}}\SpecialStringTok{ }\SpecialCharTok{\{}\DecValTok{2}\OperatorTok{\^{}}\NormalTok{dim}\SpecialCharTok{\}}\SpecialStringTok{"}\NormalTok{)}
%    \ControlFlowTok{except} \PreprocessorTok{Exception} \ImportTok{as}\NormalTok{ e:}
%        \BuiltInTok{print}\NormalTok{(}\SpecialStringTok{f"  }\SpecialCharTok{\{}\NormalTok{label}\SpecialCharTok{:\textless{}30\}}\SpecialStringTok{ error: }\SpecialCharTok{\{}\NormalTok{e}\SpecialCharTok{\}}\SpecialStringTok{"}\NormalTok{)}
%
%\BuiltInTok{print}\NormalTok{()}
%
%\CommentTok{\# ============================================================}
%\CommentTok{\# 13. LFSR simulation for systematic encoding}
%\CommentTok{\# ============================================================}
%\BuiltInTok{print}\NormalTok{(}\StringTok{"="} \OperatorTok{*} \DecValTok{70}\NormalTok{)}
%\BuiltInTok{print}\NormalTok{(}\StringTok{"13. LFSR simulation: systematic encoding with g(X)=1+X+X\^{}3"}\NormalTok{)}
%\BuiltInTok{print}\NormalTok{(}\StringTok{"="} \OperatorTok{*} \DecValTok{70}\NormalTok{)}
%
%\KeywordTok{def}\NormalTok{ lfsr\_systematic\_encode(message\_bits, g\_coeffs, n, k):}
%    \CommentTok{"""}
%\CommentTok{    Simulate systematic encoding using an LFSR.}
%\CommentTok{    message\_bits: list of k message bits [m\_0, m\_1, ..., m\_\{k{-}1\}]}
%\CommentTok{    g\_coeffs: coefficients [g\_0, g\_1, ..., g\_\{n{-}k{-}1\}]}
%\CommentTok{    Bits are fed in from highest degree to lowest.}
%\CommentTok{    """}
%\NormalTok{    r }\OperatorTok{=}\NormalTok{ n }\OperatorTok{{-}}\NormalTok{ k}
%\NormalTok{    reg }\OperatorTok{=}\NormalTok{ [F(}\DecValTok{0}\NormalTok{)] }\OperatorTok{*}\NormalTok{ r}
%
%    \BuiltInTok{print}\NormalTok{(}\SpecialStringTok{f"  }\SpecialCharTok{\{}\StringTok{\textquotesingle{}clock\textquotesingle{}}\SpecialCharTok{:\textless{}6\}}\SpecialStringTok{ }\SpecialCharTok{\{}\StringTok{\textquotesingle{}input\textquotesingle{}}\SpecialCharTok{:\textless{}8\}}\SpecialStringTok{ }\SpecialCharTok{\{}\StringTok{\textquotesingle{}feedback\textquotesingle{}}\SpecialCharTok{:\textless{}10\}}\SpecialStringTok{ }\SpecialCharTok{\{}\StringTok{\textquotesingle{}registers\textquotesingle{}}\SpecialCharTok{\}}\SpecialStringTok{"}\NormalTok{)}
%    \BuiltInTok{print}\NormalTok{(}\SpecialStringTok{f"  }\SpecialCharTok{\{}\StringTok{\textquotesingle{}init\textquotesingle{}}\SpecialCharTok{:\textless{}6\}}\SpecialStringTok{ }\SpecialCharTok{\{}\StringTok{\textquotesingle{}{-}\textquotesingle{}}\SpecialCharTok{:\textless{}8\}}\SpecialStringTok{ }\SpecialCharTok{\{}\StringTok{\textquotesingle{}{-}\textquotesingle{}}\SpecialCharTok{:\textless{}10\}}\SpecialStringTok{ }\SpecialCharTok{\{}\NormalTok{reg}\SpecialCharTok{\}}\SpecialStringTok{"}\NormalTok{)}
%
%    \ControlFlowTok{for}\NormalTok{ clk }\KeywordTok{in} \BuiltInTok{range}\NormalTok{(k):}
%\NormalTok{        input\_bit }\OperatorTok{=}\NormalTok{ F(message\_bits[k }\OperatorTok{{-}} \DecValTok{1} \OperatorTok{{-}}\NormalTok{ clk])}
%\NormalTok{        feedback }\OperatorTok{=}\NormalTok{ input\_bit }\OperatorTok{+}\NormalTok{ reg[r }\OperatorTok{{-}} \DecValTok{1}\NormalTok{]}
%
%\NormalTok{        new\_reg }\OperatorTok{=}\NormalTok{ [F(}\DecValTok{0}\NormalTok{)] }\OperatorTok{*}\NormalTok{ r}
%\NormalTok{        new\_reg[}\DecValTok{0}\NormalTok{] }\OperatorTok{=}\NormalTok{ feedback }\OperatorTok{*}\NormalTok{ F(g\_coeffs[}\DecValTok{0}\NormalTok{])}
%        \ControlFlowTok{for}\NormalTok{ j }\KeywordTok{in} \BuiltInTok{range}\NormalTok{(}\DecValTok{1}\NormalTok{, r):}
%\NormalTok{            new\_reg[j] }\OperatorTok{=}\NormalTok{ reg[j }\OperatorTok{{-}} \DecValTok{1}\NormalTok{] }\OperatorTok{+}\NormalTok{ feedback }\OperatorTok{*}\NormalTok{ F(g\_coeffs[j])}
%
%\NormalTok{        reg }\OperatorTok{=}\NormalTok{ new\_reg}
%        \BuiltInTok{print}\NormalTok{(}\SpecialStringTok{f"  }\SpecialCharTok{\{}\NormalTok{clk}\OperatorTok{+}\DecValTok{1}\SpecialCharTok{:\textless{}6\}}\SpecialStringTok{ }\SpecialCharTok{\{}\BuiltInTok{int}\NormalTok{(input\_bit)}\SpecialCharTok{:\textless{}8\}}\SpecialStringTok{ }\SpecialCharTok{\{}\BuiltInTok{int}\NormalTok{(feedback)}\SpecialCharTok{:\textless{}10\}}\SpecialStringTok{ }\SpecialCharTok{\{}\NormalTok{[}\BuiltInTok{int}\NormalTok{(x) }\ControlFlowTok{for}\NormalTok{ x }\KeywordTok{in}\NormalTok{ reg]}\SpecialCharTok{\}}\SpecialStringTok{"}\NormalTok{)}
%
%    \ControlFlowTok{return}\NormalTok{ [}\BuiltInTok{int}\NormalTok{(x) }\ControlFlowTok{for}\NormalTok{ x }\KeywordTok{in}\NormalTok{ reg]}
%
%\NormalTok{g\_coeffs }\OperatorTok{=}\NormalTok{ [}\DecValTok{1}\NormalTok{, }\DecValTok{1}\NormalTok{, }\DecValTok{0}\NormalTok{]}
%\NormalTok{message }\OperatorTok{=}\NormalTok{ [}\DecValTok{1}\NormalTok{, }\DecValTok{0}\NormalTok{, }\DecValTok{1}\NormalTok{, }\DecValTok{1}\NormalTok{]}
%
%\BuiltInTok{print}\NormalTok{(}\SpecialStringTok{f"Message: }\SpecialCharTok{\{}\NormalTok{message}\SpecialCharTok{\}}\SpecialStringTok{"}\NormalTok{)}
%\BuiltInTok{print}\NormalTok{(}\SpecialStringTok{f"g(X) = 1 + X + X\^{}3, feedback coefficients: g\_0=}\SpecialCharTok{\{}\NormalTok{g\_coeffs[}\DecValTok{0}\NormalTok{]}\SpecialCharTok{\}}\SpecialStringTok{, g\_1=}\SpecialCharTok{\{}\NormalTok{g\_coeffs[}\DecValTok{1}\NormalTok{]}\SpecialCharTok{\}}\SpecialStringTok{, g\_2=}\SpecialCharTok{\{}\NormalTok{g\_coeffs[}\DecValTok{2}\NormalTok{]}\SpecialCharTok{\}}\SpecialStringTok{"}\NormalTok{)}
%\BuiltInTok{print}\NormalTok{()}
%
%\NormalTok{parity\_lfsr }\OperatorTok{=}\NormalTok{ lfsr\_systematic\_encode(message, g\_coeffs, n, k)}
%\BuiltInTok{print}\NormalTok{(}\SpecialStringTok{f"}\CharTok{\textbackslash{}n}\SpecialStringTok{LFSR output (register state): }\SpecialCharTok{\{}\NormalTok{parity\_lfsr}\SpecialCharTok{\}}\SpecialStringTok{"}\NormalTok{)}
%
%\NormalTok{m\_poly\_check }\OperatorTok{=} \DecValTok{1} \OperatorTok{+}\NormalTok{ X}\OperatorTok{\^{}}\DecValTok{2} \OperatorTok{+}\NormalTok{ X}\OperatorTok{\^{}}\DecValTok{3}
%\NormalTok{\_, p\_check }\OperatorTok{=}\NormalTok{ (X}\OperatorTok{\^{}}\DecValTok{3} \OperatorTok{*}\NormalTok{ m\_poly\_check).quo\_rem(g)}
%\NormalTok{p\_check\_vec }\OperatorTok{=}\NormalTok{ [}\BuiltInTok{int}\NormalTok{(p\_check[i]) }\ControlFlowTok{for}\NormalTok{ i }\KeywordTok{in} \BuiltInTok{range}\NormalTok{(}\DecValTok{3}\NormalTok{)]}
%\BuiltInTok{print}\NormalTok{(}\SpecialStringTok{f"Parity from polynomial division: }\SpecialCharTok{\{}\NormalTok{p\_check\_vec}\SpecialCharTok{\}}\SpecialStringTok{"}\NormalTok{)}
%\BuiltInTok{print}\NormalTok{()}
%
%\CommentTok{\# ============================================================}
%\CommentTok{\# 14. Example of a (7,3) cyclic code (Example 5.3)}
%\CommentTok{\# ============================================================}
%\BuiltInTok{print}\NormalTok{(}\StringTok{"="} \OperatorTok{*} \DecValTok{70}\NormalTok{)}
%\BuiltInTok{print}\NormalTok{(}\StringTok{"14. (7,3) cyclic code: g(X) = 1 + X\^{}2 + X\^{}3 + X\^{}4"}\NormalTok{)}
%\BuiltInTok{print}\NormalTok{(}\StringTok{"="} \OperatorTok{*} \DecValTok{70}\NormalTok{)}
%
%\NormalTok{g73 }\OperatorTok{=} \DecValTok{1} \OperatorTok{+}\NormalTok{ X}\OperatorTok{\^{}}\DecValTok{2} \OperatorTok{+}\NormalTok{ X}\OperatorTok{\^{}}\DecValTok{3} \OperatorTok{+}\NormalTok{ X}\OperatorTok{\^{}}\DecValTok{4}
%\NormalTok{C73 }\OperatorTok{=}\NormalTok{ codes.CyclicCode(generator\_pol}\OperatorTok{=}\NormalTok{g73, length}\OperatorTok{=}\DecValTok{7}\NormalTok{)}
%\BuiltInTok{print}\NormalTok{(}\SpecialStringTok{f"Code parameters: [}\SpecialCharTok{\{}\NormalTok{C73}\SpecialCharTok{.}\NormalTok{length()}\SpecialCharTok{\}}\SpecialStringTok{, }\SpecialCharTok{\{}\NormalTok{C73}\SpecialCharTok{.}\NormalTok{dimension()}\SpecialCharTok{\}}\SpecialStringTok{, }\SpecialCharTok{\{}\NormalTok{C73}\SpecialCharTok{.}\NormalTok{minimum\_distance()}\SpecialCharTok{\}}\SpecialStringTok{]"}\NormalTok{)}
%\BuiltInTok{print}\NormalTok{(}\SpecialStringTok{f"Generator polynomial: g(X) = }\SpecialCharTok{\{}\NormalTok{g73}\SpecialCharTok{\}}\SpecialStringTok{"}\NormalTok{)}
%\BuiltInTok{print}\NormalTok{()}
%
%\NormalTok{G73 }\OperatorTok{=}\NormalTok{ Matrix(F, }\DecValTok{3}\NormalTok{, }\DecValTok{7}\NormalTok{)}
%\ControlFlowTok{for}\NormalTok{ i }\KeywordTok{in} \BuiltInTok{range}\NormalTok{(}\DecValTok{3}\NormalTok{):}
%\NormalTok{    poly }\OperatorTok{=}\NormalTok{ (X}\OperatorTok{\^{}}\NormalTok{i }\OperatorTok{*}\NormalTok{ g73) }\OperatorTok{\%}\NormalTok{ (X}\OperatorTok{\^{}}\DecValTok{7} \OperatorTok{+} \DecValTok{1}\NormalTok{)}
%    \ControlFlowTok{for}\NormalTok{ j }\KeywordTok{in} \BuiltInTok{range}\NormalTok{(}\DecValTok{7}\NormalTok{):}
%\NormalTok{        G73[i, j] }\OperatorTok{=}\NormalTok{ poly[j]}
%
%\BuiltInTok{print}\NormalTok{(}\StringTok{"Generator matrix G:"}\NormalTok{)}
%\BuiltInTok{print}\NormalTok{(G73)}
%\BuiltInTok{print}\NormalTok{()}
%
%\NormalTok{m73 }\OperatorTok{=}\NormalTok{ vector(F, [}\DecValTok{1}\NormalTok{, }\DecValTok{1}\NormalTok{, }\DecValTok{0}\NormalTok{])}
%\NormalTok{c73 }\OperatorTok{=}\NormalTok{ m73 }\OperatorTok{*}\NormalTok{ G73}
%\BuiltInTok{print}\NormalTok{(}\SpecialStringTok{f"m = (1,1,0) {-}\textgreater{} c = }\SpecialCharTok{\{}\BuiltInTok{list}\NormalTok{(c73)}\SpecialCharTok{\}}\SpecialStringTok{"}\NormalTok{)}
%\BuiltInTok{print}\NormalTok{()}
%
%\CommentTok{\# ============================================================}
%\CommentTok{\# 15. Comparison with SageMath\textquotesingle{}s built{-}in CyclicCode}
%\CommentTok{\# ============================================================}
%\BuiltInTok{print}\NormalTok{(}\StringTok{"="} \OperatorTok{*} \DecValTok{70}\NormalTok{)}
%\BuiltInTok{print}\NormalTok{(}\StringTok{"15. Comparing manual computation with SageMath\textquotesingle{}s built{-}in CyclicCode"}\NormalTok{)}
%\BuiltInTok{print}\NormalTok{(}\StringTok{"="} \OperatorTok{*} \DecValTok{70}\NormalTok{)}
%
%\NormalTok{C\_builtin }\OperatorTok{=}\NormalTok{ codes.CyclicCode(generator\_pol}\OperatorTok{=}\NormalTok{g, length}\OperatorTok{=}\DecValTok{7}\NormalTok{)}
%\NormalTok{G\_builtin }\OperatorTok{=}\NormalTok{ C\_builtin.systematic\_generator\_matrix()}
%\BuiltInTok{print}\NormalTok{(}\StringTok{"Systematic generator matrix from SageMath:"}\NormalTok{)}
%\BuiltInTok{print}\NormalTok{(G\_builtin)}
%\BuiltInTok{print}\NormalTok{()}
%
%\NormalTok{match\_count }\OperatorTok{=} \DecValTok{0}
%\NormalTok{total }\OperatorTok{=} \DecValTok{2}\OperatorTok{\^{}}\NormalTok{k}
%\ControlFlowTok{for}\NormalTok{ m\_vec }\KeywordTok{in}\NormalTok{ VectorSpace(F, k):}
%\NormalTok{    c\_builtin }\OperatorTok{=}\NormalTok{ m\_vec }\OperatorTok{*}\NormalTok{ G\_builtin}
%\NormalTok{    m\_poly }\OperatorTok{=} \BuiltInTok{sum}\NormalTok{(F(m\_vec[i]) }\OperatorTok{*}\NormalTok{ X}\OperatorTok{\^{}}\NormalTok{i }\ControlFlowTok{for}\NormalTok{ i }\KeywordTok{in} \BuiltInTok{range}\NormalTok{(k))}
%\NormalTok{    \_, \_, c\_manual }\OperatorTok{=}\NormalTok{ systematic\_encode(m\_poly, g, n, k)}
%
%    \ControlFlowTok{if} \BuiltInTok{list}\NormalTok{(c\_builtin) }\OperatorTok{==} \BuiltInTok{list}\NormalTok{(c\_manual):}
%\NormalTok{        match\_count }\OperatorTok{+=} \DecValTok{1}
%
%\BuiltInTok{print}\NormalTok{(}\SpecialStringTok{f"Matching codewords: }\SpecialCharTok{\{}\NormalTok{match\_count}\SpecialCharTok{\}}\SpecialStringTok{/}\SpecialCharTok{\{}\NormalTok{total}\SpecialCharTok{\}}\SpecialStringTok{"}\NormalTok{)}
%\BuiltInTok{print}\NormalTok{(}\SpecialStringTok{f"Manual encoding and SageMath built{-}in result }\SpecialCharTok{\{}\StringTok{\textquotesingle{}match perfectly ✓\textquotesingle{}} \ControlFlowTok{if}\NormalTok{ match\_count }\OperatorTok{==}\NormalTok{ total }\ControlFlowTok{else} \StringTok{\textquotesingle{}do not fully match ✗\textquotesingle{}}\SpecialCharTok{\}}\SpecialStringTok{"}\NormalTok{)}
%\BuiltInTok{print}\NormalTok{()}
%
%\CommentTok{\# ============================================================}
%\CommentTok{\# 16. Structure of a circulant block (Example 16.1)}
%\CommentTok{\# ============================================================}
%\BuiltInTok{print}\NormalTok{(}\StringTok{"="} \OperatorTok{*} \DecValTok{70}\NormalTok{)}
%\BuiltInTok{print}\NormalTok{(}\StringTok{"16. Circulant{-}matrix block structure"}\NormalTok{)}
%\BuiltInTok{print}\NormalTok{(}\StringTok{"="} \OperatorTok{*} \DecValTok{70}\NormalTok{)}
%
%\KeywordTok{def}\NormalTok{ circulant\_matrix(first\_row, field}\OperatorTok{=}\NormalTok{F):}
%    \CommentTok{"""Construct a circulant matrix from its first row."""}
%\NormalTok{    m }\OperatorTok{=} \BuiltInTok{len}\NormalTok{(first\_row)}
%\NormalTok{    M }\OperatorTok{=}\NormalTok{ Matrix(field, m, m)}
%    \ControlFlowTok{for}\NormalTok{ i }\KeywordTok{in} \BuiltInTok{range}\NormalTok{(m):}
%        \ControlFlowTok{for}\NormalTok{ j }\KeywordTok{in} \BuiltInTok{range}\NormalTok{(m):}
%\NormalTok{            M[i, j] }\OperatorTok{=}\NormalTok{ first\_row[(j }\OperatorTok{{-}}\NormalTok{ i) }\OperatorTok{\%}\NormalTok{ m]}
%    \ControlFlowTok{return}\NormalTok{ M}
%
%\NormalTok{first\_row }\OperatorTok{=}\NormalTok{ [F(}\DecValTok{1}\NormalTok{), F(}\DecValTok{0}\NormalTok{), F(}\DecValTok{1}\NormalTok{), F(}\DecValTok{1}\NormalTok{)]}
%\NormalTok{C\_circ }\OperatorTok{=}\NormalTok{ circulant\_matrix(first\_row)}
%\BuiltInTok{print}\NormalTok{(}\SpecialStringTok{f"First row: }\SpecialCharTok{\{}\NormalTok{first\_row}\SpecialCharTok{\}}\SpecialStringTok{"}\NormalTok{)}
%\BuiltInTok{print}\NormalTok{(}\StringTok{"Circulant matrix:"}\NormalTok{)}
%\BuiltInTok{print}\NormalTok{(C\_circ)}
%\BuiltInTok{print}\NormalTok{()}
%
%\NormalTok{circ\_ok }\OperatorTok{=} \VariableTok{True}
%\ControlFlowTok{for}\NormalTok{ i }\KeywordTok{in} \BuiltInTok{range}\NormalTok{(}\DecValTok{1}\NormalTok{, }\DecValTok{4}\NormalTok{):}
%\NormalTok{    prev\_row }\OperatorTok{=} \BuiltInTok{list}\NormalTok{(C\_circ[i }\OperatorTok{{-}} \DecValTok{1}\NormalTok{])}
%\NormalTok{    expected }\OperatorTok{=}\NormalTok{ [prev\_row[}\OperatorTok{{-}}\DecValTok{1}\NormalTok{]] }\OperatorTok{+}\NormalTok{ prev\_row[:}\OperatorTok{{-}}\DecValTok{1}\NormalTok{]}
%\NormalTok{    actual }\OperatorTok{=} \BuiltInTok{list}\NormalTok{(C\_circ[i])}
%    \ControlFlowTok{if}\NormalTok{ actual }\OperatorTok{!=}\NormalTok{ expected:}
%\NormalTok{        circ\_ok }\OperatorTok{=} \VariableTok{False}
%        \BuiltInTok{print}\NormalTok{(}\SpecialStringTok{f"  ✗ row }\SpecialCharTok{\{}\NormalTok{i}\SpecialCharTok{\}}\SpecialStringTok{: expected }\SpecialCharTok{\{}\NormalTok{expected}\SpecialCharTok{\}}\SpecialStringTok{, got }\SpecialCharTok{\{}\NormalTok{actual}\SpecialCharTok{\}}\SpecialStringTok{"}\NormalTok{)}
%
%\ControlFlowTok{if}\NormalTok{ circ\_ok:}
%    \BuiltInTok{print}\NormalTok{(}\StringTok{"  ✓ Every row is a cyclic right shift of the previous row."}\NormalTok{)}
%\BuiltInTok{print}\NormalTok{()}
%
%\BuiltInTok{print}\NormalTok{(}\StringTok{"="} \OperatorTok{*} \DecValTok{70}\NormalTok{)}
%\BuiltInTok{print}\NormalTok{(}\StringTok{"All checks completed!"}\NormalTok{)}
%\BuiltInTok{print}\NormalTok{(}\StringTok{"="} \OperatorTok{*} \DecValTok{70}\NormalTok{)}
%\end{Highlighting}
%\end{Shaded}


\backmatter
\chapter*{References}\addcontentsline{toc}{chapter}{References}\label{sec:references}

{[}1{]} C. E. Shannon, ``A Mathematical Theory of Communication,''
\emph{Bell System Technical Journal}, vol.~27, no. 3, pp.~379-423, 1948.

{[}2{]} R. W. Hamming, ``Error Detecting and Error Correcting Codes,''
\emph{Bell System Technical Journal}, vol.~29, no. 2, pp.~147-160, 1950.

{[}3{]} D. E. Muller, ``Application of Boolean Algebra to Switching
Circuit Design and to Error Detection,'' \emph{IRE Transactions on
Electronic Computers}, vol.~EC-3, no. 3, pp.~6-12, 1954.

{[}4{]} I. S. Reed, ``A Class of Multiple-Error-Correcting Codes and the
Decoding Scheme,'' \emph{IRE Transactions on Information Theory},
vol.~4, no. 4, pp.~38-49, 1954.

{[}5{]} E. Prange, ``Cyclic Error-Correcting Codes in Two Symbols,''
\emph{AFCRC-TN-57-103}, Air Force Cambridge Research Center, 1957.

{[}6{]} A. Hocquenghem, ``Codes correcteurs d'erreurs,''
\emph{Chiffres}, vol.~2, pp.~147-156, 1959.

{[}7{]} R. C. Bose and D. K. Ray-Chaudhuri, ``On a Class of Error
Correcting Binary Group Codes,'' \emph{Information and Control}, vol.~3,
no. 1, pp.~68-79, 1960.

{[}8{]} I. S. Reed and G. Solomon, ``Polynomial Codes Over Certain
Finite Fields,'' \emph{Journal of the Society for Industrial and Applied
Mathematics}, vol.~8, no. 2, pp.~300-304, 1960.

{[}9{]} W. W. Peterson, \emph{Error-Correcting Codes}, MIT Press, 1961.

{[}10{]} E. R. Berlekamp, \emph{Algebraic Coding Theory}, McGraw-Hill,
1968.

{[}11{]} J. L. Massey, ``Shift-Register Synthesis and BCH Decoding,''
\emph{IEEE Transactions on Information Theory}, vol.~15, no. 1,
pp.~122-127, 1969.

{[}12{]} V. D. Goppa, ``Geometry and Codes,'' \emph{Mathematics and its
Applications}, vol.~24, Kluwer, 1988.

{[}13{]} R. J. McEliece, ``A Public-Key Cryptosystem Based on Algebraic
Coding Theory,'' \emph{DSN Progress Report}, Jet Propulsion Laboratory,
1978.

{[}14{]} C. Berrou, A. Glavieux, and P. Thitimajshima, ``Near Shannon
Limit Error-Correcting Coding and Decoding: Turbo Codes,'' in
\emph{Proceedings of ICC '93}, 1993.

{[}15{]} D. J. C. MacKay and R. M. Neal, ``Near Shannon Limit
Performance of Low Density Parity Check Codes,'' \emph{Electronics
Letters}, vol.~32, no. 18, pp.~1645-1646, 1996.


\end{document}
